% !TEX program = xelatex
\documentclass[11pt,a4paper]{article}
\usepackage[a4paper,margin=2.3cm]{geometry}
\usepackage{fontspec}
\usepackage{amsmath,amssymb,amsthm}
\usepackage{xcolor}
\usepackage{booktabs,tabularx,longtable,array}
\usepackage{enumitem}
\usepackage{fvextra}
\usepackage{newunicodechar}
\usepackage[most]{tcolorbox}
\usepackage{url}
\usepackage[colorlinks=true,linkcolor=blue!50!black,urlcolor=blue!55!black,citecolor=blue!50!black]{hyperref}
\usepackage{fancyhdr}
\usepackage{microtype}

% ---------- fonts: a monospace face with wide Unicode coverage for Lean source ----------
\setmonofont{DejaVuSansMono.ttf}[Scale=0.80,BoldFont=DejaVuSansMono-Bold.ttf,ItalicFont=DejaVuSansMono-Oblique.ttf,BoldItalicFont=DejaVuSansMono-BoldOblique.ttf]
\newfontfamily\mathfb{STIX Two Math}[Scale=0.88]
\newunicodechar{𝒞}{{\mathfb 𝒞}}
\newunicodechar{𝓞}{{\mathfb 𝓞}}
\newunicodechar{𝒜}{{\mathfb 𝒜}}
\newunicodechar{𝒮}{{\mathfb 𝒮}}
\newunicodechar{𝒟}{{\mathfb 𝒟}}
\newunicodechar{ℰ}{{\mathfb ℰ}}

% ---------- macros ----------
\makeatletter\g@addto@macro{\UrlBreaks}{\do\_\do\.}\makeatother
\Urlmuskip=0mu plus 2mu
\DeclareUrlCommand\lean{\urlstyle{tt}}
\DeclareUrlCommand\path{\urlstyle{tt}}
\newcommand{\leanu}[1]{\texttt{#1}}% identifiers containing non-ASCII characters or spaces
\newcommand{\mtt}[1]{\text{\texttt{#1}}}
\newcommand{\N}{\textup{[N]}}% the paper under review
\newcommand{\OK}{\textcolor{green!40!black}{\textbf{correct}}}
\newcommand{\PARTLY}{\textcolor{orange!70!black}{\textbf{partly}}}
\newcommand{\NO}{\textcolor{red!70!black}{\textbf{incorrect}}}
\definecolor{boxgray}{gray}{0.955}
\newtcolorbox{verdictbox}[1][]{enhanced,breakable,colback=blue!3,colframe=blue!45!black,boxrule=0.8pt,arc=2pt,left=7pt,right=7pt,top=6pt,bottom=6pt,#1}
\newtcolorbox{notebox}[1][]{enhanced,breakable,colback=boxgray,colframe=gray!55,boxrule=0.5pt,arc=2pt,left=7pt,right=7pt,top=5pt,bottom=5pt,#1}

\makeatletter\newcommand{\tableinput}[1]{\@@input #1 }\makeatother
\expandafter\def\csname ex@anygeo@first\endcsname{14}
\expandafter\def\csname ex@anygeo@last\endcsname{19}
\expandafter\def\csname ex@anygeo@file\endcsname{\path{src/Development/AbelianVarietyTheory/Definitions/AutomorphicMembership/OverIx.lean}}
\expandafter\def\csname ex@oat@first\endcsname{44}
\expandafter\def\csname ex@oat@last\endcsname{49}
\expandafter\def\csname ex@oat@file\endcsname{\path{src/Development/AbelianVarietyTheory/Definitions/AutomorphicMembership/OverIx.lean}}
\expandafter\def\csname ex@geocl@first\endcsname{21}
\expandafter\def\csname ex@geocl@last\endcsname{28}
\expandafter\def\csname ex@geocl@file\endcsname{\path{src/Development/AbelianVarietyTheory/Definitions/AutomorphicMembership/NotGeoClFalse.lean}}
\expandafter\def\csname ex@main@first\endcsname{1}
\expandafter\def\csname ex@main@last\endcsname{12}
\expandafter\def\csname ex@main@file\endcsname{\path{src/Main.lean}}
\expandafter\def\csname ex@maintheorems@first\endcsname{7}
\expandafter\def\csname ex@maintheorems@last\endcsname{9}
\expandafter\def\csname ex@maintheorems@file\endcsname{\path{src/Supplements/Kernels/MainTheorems.lean}}
\expandafter\def\csname ex@siegel_a@first\endcsname{59}
\expandafter\def\csname ex@siegel_a@last\endcsname{67}
\expandafter\def\csname ex@siegel_a@file\endcsname{\path{src/Development/AutomorphicArithmetic/ThetaCorrespondence/ThetaPairLifts/ThirdPlaceSeries/SiegelEisensteinSecondTermHypotheses.lean}}
\expandafter\def\csname ex@zfa@first\endcsname{12}
\expandafter\def\csname ex@zfa@last\endcsname{18}
\expandafter\def\csname ex@zfa@file\endcsname{\path{src/Development/AutomorphicArithmetic/ThetaCorrespondence/ThetaPairLifts/ThirdPlaceSeries/TopPlug.lean}}
\expandafter\def\csname ex@shapeof@first\endcsname{31}
\expandafter\def\csname ex@shapeof@last\endcsname{32}
\expandafter\def\csname ex@shapeof@file\endcsname{\path{src/Development/AutomorphicArithmetic/ThetaCorrespondence/ThetaPairLifts/ThirdPlaceSeries/ShapeOf.lean}}
\expandafter\def\csname ex@topshape@first\endcsname{33}
\expandafter\def\csname ex@topshape@last\endcsname{33}
\expandafter\def\csname ex@topshape@file\endcsname{\path{src/Development/AutomorphicArithmetic/ThetaCorrespondence/ThetaPairLifts/ThirdPlaceSeries/TopShape.lean}}
\expandafter\def\csname ex@topg@first\endcsname{12}
\expandafter\def\csname ex@topg@last\endcsname{14}
\expandafter\def\csname ex@topg@file\endcsname{\path{src/Development/AutomorphicArithmetic/CharacterBoxes/ZeroFreeAllOfTopG.lean}}
\expandafter\def\csname ex@zfatopg@first\endcsname{26}
\expandafter\def\csname ex@zfatopg@last\endcsname{34}
\expandafter\def\csname ex@zfatopg@file\endcsname{\path{src/Development/AutomorphicArithmetic/CharacterBoxes/ZeroFreeAllOfTopG.lean}}
\expandafter\def\csname ex@params@first\endcsname{19}
\expandafter\def\csname ex@params@last\endcsname{25}
\expandafter\def\csname ex@params@file\endcsname{\path{src/Development/AutomorphicArithmetic/ThetaCorrespondence/ThetaPairLifts/ThirdPlaceSeries/CoreAltAltBis.lean}}
\expandafter\def\csname ex@mono_sig@first\endcsname{28}
\expandafter\def\csname ex@mono_sig@last\endcsname{28}
\expandafter\def\csname ex@mono_sig@file\endcsname{\path{src/Development/AutomorphicArithmetic/ThetaCorrespondence/ThetaPairLifts/ThirdPlaceSeries/OpenAndLinkedOfStrongerExtension.lean}}
\expandafter\def\csname ex@hypsalt@first\endcsname{22}
\expandafter\def\csname ex@hypsalt@last\endcsname{27}
\expandafter\def\csname ex@hypsalt@file\endcsname{\path{src/Development/AutomorphicArithmetic/ThetaCorrespondence/ThetaPairLifts/ThirdPlaceSeries/Hyps.lean}}
\expandafter\def\csname ex@divexprov@first\endcsname{204}
\expandafter\def\csname ex@divexprov@last\endcsname{209}
\expandafter\def\csname ex@divexprov@file\endcsname{\path{src/Development/AbelianVarietyTheory/Proofs/IwasawaDecomposition/TupPinsOfMuIsW.lean}}
\expandafter\def\csname ex@divexgeocl@first\endcsname{39}
\expandafter\def\csname ex@divexgeocl@last\endcsname{49}
\expandafter\def\csname ex@divexgeocl@file\endcsname{\path{src/Development/AbelianVarietyTheory/Proofs/IwasawaDecomposition/PublishedInputsImplyConclusion.lean}}
\expandafter\def\csname ex@geodegen@first\endcsname{241}
\expandafter\def\csname ex@geodegen@last\endcsname{254}
\expandafter\def\csname ex@geodegen@file\endcsname{\path{src/Formalism/LocalToGlobal/LatticeGeometry/Inhabited.lean}}
\expandafter\def\csname ex@existsdivt@first\endcsname{188}
\expandafter\def\csname ex@existsdivt@last\endcsname{191}
\expandafter\def\csname ex@existsdivt@file\endcsname{\path{src/Development/AbelianVarietyTheory/Proofs/IwasawaDecomposition/TupPinsOfMuIsW.lean}}
\expandafter\def\csname ex@heighteq@first\endcsname{379}
\expandafter\def\csname ex@heighteq@last\endcsname{400}
\expandafter\def\csname ex@heighteq@file\endcsname{\path{src/Development/AbelianVarietyTheory/Definitions/GlobalTwisting/SetUnipotentTranslate.lean}}
\expandafter\def\csname ex@hexpl6@first\endcsname{25}
\expandafter\def\csname ex@hexpl6@last\endcsname{27}
\expandafter\def\csname ex@hexpl6@file\endcsname{\path{src/Formalism/LocalRingCardinality/KernelPlaceMeanBound.lean}}
\expandafter\def\csname ex@wordfacts@first\endcsname{128}
\expandafter\def\csname ex@wordfacts@last\endcsname{133}
\expandafter\def\csname ex@wordfacts@file\endcsname{\path{src/Formalism/LocalToGlobal/Divisibility/RecBridge.lean}}
\expandafter\def\csname ex@wordfactsof@first\endcsname{135}
\expandafter\def\csname ex@wordfactsof@last\endcsname{153}
\expandafter\def\csname ex@wordfactsof@file\endcsname{\path{src/Formalism/LocalToGlobal/Divisibility/RecBridge.lean}}
\expandafter\def\csname ex@divclauses@first\endcsname{111}
\expandafter\def\csname ex@divclauses@last\endcsname{121}
\expandafter\def\csname ex@divclauses@file\endcsname{\path{src/Formalism/LocalToGlobal/Divisibility/ConstantGeTwentyEight.lean}}
\expandafter\def\csname ex@divclausesof@first\endcsname{184}
\expandafter\def\csname ex@divclausesof@last\endcsname{209}
\expandafter\def\csname ex@divclausesof@file\endcsname{\path{src/Formalism/LocalToGlobal/Divisibility/ConstantGeTwentyEight.lean}}
\expandafter\def\csname ex@divclausesof2@first\endcsname{211}
\expandafter\def\csname ex@divclausesof2@last\endcsname{212}
\expandafter\def\csname ex@divclausesof2@file\endcsname{\path{src/Formalism/LocalToGlobal/Divisibility/ConstantGeTwentyEight.lean}}
\expandafter\def\csname ex@divclausesof3@first\endcsname{245}
\expandafter\def\csname ex@divclausesof3@last\endcsname{245}
\expandafter\def\csname ex@divclausesof3@file\endcsname{\path{src/Formalism/LocalToGlobal/Divisibility/ConstantGeTwentyEight.lean}}
\expandafter\def\csname ex@kcopy@first\endcsname{31}
\expandafter\def\csname ex@kcopy@last\endcsname{42}
\expandafter\def\csname ex@kcopy@file\endcsname{\path{src/Formalism/LocalToGlobal/NumberFieldFamilies/OwnStruct.lean}}

\newcommand{\excerpt}[2][]{%
  \par\smallskip\noindent{\scriptsize\color{black!62}\csname ex@#2@file\endcsname,\ \ lines \csname ex@#2@first\endcsname--\csname ex@#2@last\endcsname}\par\vspace{-0.55em}%
  \VerbatimInput[numbers=left,firstnumber=\csname ex@#2@first\endcsname,breaklines=true,breakanywhere=true,fontsize=\scriptsize,frame=single,rulecolor=\color{gray!50},numbersep=4pt,xleftmargin=2.1em,framesep=3pt,#1]{excerpts/#2.txt}\vspace{-0.3em}}
\fvset{breaklines=true,breakanywhere=true}
% Generated by make_tables.py from evidence/depgraph/query/depquery.json. Do not edit by hand.
\newcommand{\Ndecl}{1,468,418}
\newcommand{\Nlines}{1,080,769,508}
\newcommand{\Nexprs}{1,065,214,335}
\newcommand{\NedgesT}{49,376,348}
\newcommand{\NedgesV}{106,908,310}
\newcommand{\Nedges}{156,284,658}
\newcommand{\NmathlibNames}{771,129}
\newcommand{\Nlib}{209,514}
\newcommand{\Nproj}{1,258,904}
\newcommand{\NprojThm}{1,075,031}
\newcommand{\NstmtMain}{4,566}
\newcommand{\NstmtMainProj}{0}
\newcommand{\NconeMain}{1,468,390}
\newcommand{\NconeProj}{1,258,881}
\newcommand{\NconeProjThm}{1,075,015}
\newcommand{\NanyGeo}{30}
\newcommand{\NgeoDeg}{52}
\newcommand{\Nflagged}{124}
\newcommand{\NflaggedThm}{104}
\newcommand{\NfreeProj}{1,258,757}
\newcommand{\NfreeProjThm}{1,074,911}
\newcommand{\NifaceDeps}{10,708}
\newcommand{\NaccDirect}{75}
\newcommand{\NaccThmProof}{20}
\newcommand{\Nacc}{18}
\newcommand{\NdivInputs}{38}
\newcommand{\NdivClausesAll}{71,597}
\newcommand{\NdivClausesProj}{1,418}
\newcommand{\NdivClausesThm}{1,004}
\newcommand{\NgreenAll}{228,101}
\newcommand{\NgreenProj}{95,675}
\newcommand{\NgreenThm}{74,917}
\newcommand{\NpackageAll}{932,347}
\newcommand{\NpackageProj}{744,434}
\newcommand{\NpackageThm}{636,383}
\newcommand{\NplaceLP}{3,856}
\newcommand{\NplaceNT}{3,828}
\newcommand{\NgeoTiedP}{3,481}
\newcommand{\NgeoOver}{n/a}
\newcommand{\NdivXdeps}{10,530}
% Text fragments that report findings of agent F (Gemini 4 Neon), kept in one place.
\newcommand{\Fnote}{Agent~F reached the same conclusion from the source: the structure that carries the vacuous field \lean{degHH_def} is never constructed anywhere in the release, and the statements containing the degree equation are off the route (\textsf{G-F}).}
\newcommand{\Frow}{; \textsf{G-F}: on the route \lean{InGeo} is established by lemmas and by projections of proved statements, no field by \lean{trivial} or \lean{rfl}}
\newcommand{\Fsix}{; \textsf{G-F}: each of the six occurs once in the source}
\newcommand{\Fsummary}{Confirmed the dead extraction lemmas and the storing of the certificate; \textbf{showed that the row-side condition} \lean{InGeo} \textbf{is established on the route by substantive proofs}, at a primary construction site that \N\ does not list;
\textbf{found that the structure carrying the vacuous field} \lean{degHH_def} \textbf{is never constructed}, and that the degree equation is off the route.}
\newcommand{\edgesHash}{82b5364c5d2b6f76a024ed95d704fa00c7ed7017c0119244072f45adb1058a85}


\theoremstyle{plain}
\newtheorem{proposition}{Proposition}
\theoremstyle{definition}
\newtheorem{definition}{Definition}
\theoremstyle{remark}
\newtheorem{remark}{Remark}

\setlist{itemsep=0.15em,topsep=0.3em,parsep=0.1em}
\setlength{\emergencystretch}{3.5em}
\pagestyle{fancy}\fancyhf{}
\fancyhead[L]{\small The Invalidity of ``The Invalidity of Alpöge's Proof of the Riemann Hypothesis''}
\fancyhead[R]{\small Z.~Zeng}
\fancyfoot[C]{\small\thepage}
\renewcommand{\headrulewidth}{0.3pt}
\setlength{\headheight}{14pt}
\hypersetup{pdftitle={The Invalidity of "The Invalidity of Alpöge's Proof of the Riemann Hypothesis"},pdfauthor={Zijian Zeng},
  pdfsubject={A kernel-level audit of the vacuous-witness objection to the Lean formalization},pdfkeywords={Riemann Hypothesis, Lean 4, formal verification, proof audit}}

\begin{document}

\begin{center}
{\LARGE\bfseries The Invalidity of\\[0.3em] ``The Invalidity of Alpöge's Proof\\[0.3em] of the Riemann Hypothesis''\par}
\vspace{0.9em}
{\large A kernel-level audit of the vacuous-witness objection to the Lean formalization\par}
\vspace{1.4em}
{\large Zijian Zeng\par}
\vspace{0.35em}
{\small Institute of Computer Science and Digital Innovation, UCSI University,\\ Kuala Lumpur, 56000, MALAYSIA\par}
\vspace{0.3em}
{\small \href{mailto:zijianzeng@foxmail.com}{\texttt{zijianzeng@foxmail.com}};\ \ \href{mailto:1002266693@ucsiuniversity.edu.my}{\texttt{1002266693@ucsiuniversity.edu.my}}\par}
\vspace{0.9em}
{\small Audit date: 23 November 2026\par}
\end{center}
\vspace{0.6em}

\begin{abstract}
\noindent
Nielsen \N\ argues that the Lean~4 formalization released by Alpöge on 20 November 2026 does not contain a proof of the Riemann Hypothesis:
on the route to the final theorem a geometry predicate is the constant \lean{True}, the geometric witness is a degenerate model,
lattice identifications hold by \lean{rfl}, and a height is installed by assignment.
We audit that argument against the published artefacts, namely the author's source archive and the 61.57\,GB kernel-level export of the proof term.
The source-level observations of \N\ are, with minor corrections, accurate, and we confirm at kernel level that the flagged objects lie in the dependency cone of \lean{riemannHypothesis}.
The inference from them to invalidity nevertheless fails.
(i)~What the kernel certifies is a closed term whose type is Mathlib's proposition \lean{RiemannHypothesis} and which uses only the three standard axioms.
Its statement depends on \NstmtMain\ declarations, all of them in Mathlib and none in the release; it is not an identity of the form $X = X$, and auxiliary definitions cannot change what it says.
(ii)~The claim that the central squeeze compares a quantity with itself and holds by reflexivity is contradicted by the code and by Section~7 of \N\ itself:
the two inequalities are proved from \NdivInputs\ hypotheses (operator-norm, valuation and covolume bounds) by a theorem whose dependency cone contains \NdivClausesThm\ theorems of the release and none of the flagged objects.
(iii)~The reply of \N\ to the ``kernel objection'' tests soundness by \emph{deleting} an input. The relevant test is \emph{replacing} it, and the consumers on the route are universally quantified over the inputs in question.
(iv)~Only \NanyGeo\ of \Ndecl\ declarations depend on the \lean{True} predicate, while the hypothesis package it is said to hollow out is established using \NpackageThm\ theorems of the release.
What \N\ does establish is weaker: the formal proof does not use the geometric interpretation that its identifiers advertise.
That concerns exposition, not validity.
We also report an independent replay of the complete export with the \texttt{con-leche} checker and a byte-identical rebuild of the Mathlib-only statement.
\end{abstract}

\section{Introduction}
\label{sec:intro}

On 20 November 2026 Levent Alpöge announced a proof of the Riemann Hypothesis, and of the generalized Riemann Hypothesis for Dirichlet $L$-functions, obtained with the AI system Claude~\cite{alpoge}.
The announcement was accompanied by a Lean~4 development of 175,187 source files and by no human-readable paper.
Within a day the Lean and Mathlib communities reported that the statements are the intended ones and that the proof checks~\cite{community}, and a verification page with the complete kernel-level export of the proof and the logs of several independent checkers was published~\cite{mit}.

Nielsen~\N\ disputes the result. The paper reads the source and reports that the part of the development named after arithmetic geometry is empty of geometric content:
a geometry predicate is defined to be \lean{True}, the geometric datum is assembled from objects that the source itself calls \emph{degenerate} and \emph{junk},
the lattice conditions on that datum hold by \lean{rfl} because the datum was defined as a copy, and the height that enters the central inequality is installed by assignment.
From this \N\ concludes that the formalization certifies only a tautology and that the release contains no proof of the Riemann Hypothesis.

This note audits that argument. It does not try to prove or to disprove the Riemann Hypothesis, and it is not a referee report on the mathematics of the release.
It asks two narrower questions that can be settled by inspecting the published artefacts:

\begin{enumerate}[label=\textbf{Q\arabic*.},leftmargin=3em]
\item Are the observations of \N\ about the code true?
\item If they are true, does the conclusion of \N\ follow?
\end{enumerate}

\begin{verdictbox}
\textbf{Findings.}
\textbf{Q1: largely yes.} We confirm the observations in the source, and we confirm at the level of the kernel export that the objects flagged by \N\ lie in the dependency cone of \lean{riemannHypothesis}
(Section~\ref{sec:facts}). A few citations in \N\ are inaccurate, and two headline descriptions (that the squeeze holds ``by reflexivity'', and that what is certified is a ``tautology'') are false as statements about the code.

\textbf{Q2: no.} The step from ``these witnesses are trivial'' to ``the theorem is unproved'' is not valid for a closed, kernel-checked theorem (Section~\ref{sec:inference}).
A trivially discharged slot in such a proof is a slot the proof does not need. The objection would have force against a theorem that still carried the weakened predicate as a hypothesis, or against one that used an axiom; the final theorem here does neither.
We looked independently for the ways in which a kernel-checked theorem can nevertheless be wrong (extra axioms, a substituted statement, an exploited kernel defect, a mismatch between source and checked artefact) and found none (Section~\ref{sec:defects}).
\end{verdictbox}

\paragraph{What is and is not claimed here.}
The claim defended is conditional and precise: \emph{given} that the published export is accepted by the Lean kernel and by independent checkers as a derivation of Mathlib's \lean{RiemannHypothesis} from the three standard axioms,
the observations of \N\ give no reason to doubt that it is one. The trust base that remains (the soundness of Lean's logic, the correctness of at least one kernel implementation, and the adequacy of Mathlib's definition of the Riemann zeta function)
is discussed in Section~\ref{sec:limits}; \N\ does not challenge any part of it.
We have not read the mathematics of the release and we express no view on why the proof works.
Section~\ref{sec:credit} states what \N\ does establish, which is of real interest for anyone who wants to understand the proof.

\paragraph{Method.}
Every source excerpt in this note is cut mechanically, with its true line numbers, from the author's original archive, whose hash we verified.
Every statement about what the proof term \emph{uses} is derived from the kernel-level export~\cite{mit} rather than from text search:
we wrote a program that extracts the complete declaration-level dependency graph of the export (\Ndecl\ declarations, \Nedges\ edges) and answer reachability questions on that graph.
The tools, their outputs, the logs of the independent checker run and the reports of eight read-only cross-check agents are in the accompanying repository~\cite{repo}.
An earlier version of this audit was written from the abstract and public summaries of \N; the present version is based on the full text.
Appendix~\ref{app:erratum} corrects an undercount in that earlier version.

\section{The artefacts}
\label{sec:objects}

Table~\ref{tab:prov} lists what was audited. Each hash was recomputed on our machine and compared with the published value.

\begin{table}[ht]
\centering\small
\renewcommand{\arraystretch}{1.2}
\begin{tabularx}{\textwidth}{@{}>{\raggedright\arraybackslash}p{3.7cm}>{\raggedright\arraybackslash}X>{\raggedright\arraybackslash}p{4.4cm}@{}}
\toprule
Object & Origin and SHA-256 (computed locally) & Check \\
\midrule
Author's source archive \path{lean.tar.gz} (239,275,554 bytes) & 14 parts from \path{alpo.ge}, concatenated; \texttt{9631b4f0…93f49c7fc} & equals the hash on the verification page~\cite{mit}; the same archive \N\ cites \\
Source tree (175,187 \path{.lean} files) & unpacked archive, compared by \texttt{diff -rq} with the mirror~\cite{mirror} at commit \texttt{c7d25b3b} & \path{src/} identical file by file; the mirror adds only the community file \path{Challenge.lean} \\
Proof export \path{solution.ndjson.zst} (7,699,933,588 bytes) & verification page; \texttt{61a8938a…a681b05381} & equals the published value \\
Unpacked \path{solution.ndjson} (61,570,595,788 bytes) & \texttt{5ea74b91…f6be244e} & equals the published value \\
Statement export \path{challenge.ndjson} (350,862,311 bytes) & verification page; \texttt{bd484ac9…bdd6ab2b9} & equals the published value, \textbf{and is byte-identical to our own rebuild from official Mathlib} (Section~\ref{sec:rebuild}) \\
Mathlib & pinned in \path{lakefile.toml} at \texttt{0df444a360ea…} & \texttt{git ls-remote}: this is the official tag \texttt{v4.33.1} \\
The paper \N\ (17 pages) & title page dated 20 November 2026; \texttt{90b7c28e…bc7ccf77} & the version reviewed here \\
\bottomrule
\end{tabularx}
\caption{Artefacts and integrity checks. Full hashes are in Appendix~\ref{app:hashes}.}
\label{tab:prov}
\end{table}

\paragraph{The export is the dependency cone.}
The export was produced by \texttt{lean4export} with the five headline theorems as targets, so it contains exactly the declarations that at least one of them uses, directly or indirectly.
A name that does not occur in the export is therefore not used by the proof, whatever the source text suggests.
The export has \Nlines\ lines and contains \Nexprs\ expression nodes and \Ndecl\ declaration records:
1,246,110 theorems, 214,241 definitions, 7,969 inductive blocks, 91 opaque constants, 4 quotient primitives and \textbf{3 axioms}, namely \lean{propext}, \lean{Quot.sound} and \lean{Classical.choice}.

\paragraph{Library and release.}
To separate the release from the library it builds on, we dumped the names of all \NmathlibNames\ constants of the environment obtained by \texttt{import Mathlib} at the pinned version.
A declaration of the export whose name is not among them was declared by the release. By this criterion the export consists of \Nlib\ library declarations and \Nproj\ declarations of the release,
of which \NprojThm\ are theorems. All 91 opaque constants and all three axioms belong to the library.

\paragraph{Tools.}
\begin{itemize}
\item \textbf{depgraph} (written for this audit; one Rust file, no dependencies). It reads the export once, stores the expression DAG, and for every declaration walks the DAG of its type and of its value,
recording which declarations are referenced. References found in the \emph{type} (the statement) are kept apart from those found in the \emph{value} (the proof or the definition body).
It also records every structure projection and every direct reference to a list of watched names such as recursors. The result is a graph with \Ndecl\ nodes, \NedgesT\ type edges and \NedgesV\ value edges.
\item \textbf{depquery} (Python). Reachability queries on that graph; see Definition~\ref{def:cones}.
\item \textbf{conescan} (written for the earlier version). A single-pass taint propagation over the same export. Its results for definitions and theorems agree with depgraph; see Appendix~\ref{app:erratum} for a limitation.
\item \textbf{con-leche}~\cite{conleche}. The independent checker of the Lean FRO, built from source and run on the complete export (Section~\ref{sec:conleche}).
\item \textbf{Official tool chain.} Lean 4.33.1 and 4.36.0 release binaries, \texttt{lean4export} at tag \texttt{v4.33.0}, the official Mathlib build cache.
\item \textbf{Cross-checks.} Eight read-only agents (Gemini~4 Neon), each given one sub-question and the source tree; four in the earlier round and four on the full text of \N.
Their reports are in the repository. A finding of an agent is used here only where we re-checked it in the source or in the graph, or else it is marked as an agent finding.
\end{itemize}

\begin{definition}\label{def:cones}
For a declaration $X$ of the export:
the \emph{dependency cone} of $X$ is the set of declarations that $X$ refers to transitively, through types and values;
the \emph{statement cone} of $X$ is what the statement of $X$ refers to: the type of $X$ and then, transitively, the type and value of every definition, inductive type or axiom reached, but only the \emph{type} of every theorem reached
(a proof that occurs inside a statement matters only through what it proves);
the \emph{dependents} of $X$ are the declarations whose dependency cone contains $X$.
\end{definition}

\section{The argument under review}
\label{sec:argument}

We first state the argument of \N\ in a form that separates what it observes from what it infers. Section numbers refer to \N.
The reconstruction was made twice, by us and independently by cross-check agent~H (Section~\ref{sec:agents}), with the same result.

\begin{notebox}
\textbf{Observations about the code}
\begin{enumerate}[label=(A\arabic*),leftmargin=3.2em]
\item The final theorem is a closed term of Mathlib's \lean{RiemannHypothesis}; the source has no \lean{axiom}, \lean{sorry}, \lean{native_decide}, \lean{unsafe} or \lean{opaque} declaration. \hfill (Abstract, \S12)
\item The proof term reaches it along the route \lean{riemannHypothesis} $\leftarrow$ \lean{rhOfAll} $\leftarrow$ \lean{allChi} $\leftarrow$ \lean{everyCharacter} $\leftarrow\cdots\leftarrow$
\lean{every_character_of_green_v_win} $\leftarrow$ \lean{zero_free_all_of_green_v_win} $\leftarrow$ \lean{zeroFreeAll_of_topG}. The theorem \lean{plugAll} is not on it. \hfill (\S1, \S7, \S11)
\item On this route the witness of the hypothesis bundle \lean{HypsAlt} is built with \texttt{geoCl anyGeometry}, where \lean{anyGeometry} is the constant \lean{True}; two further ``shape'' inputs are supplied by \lean{trivial}. \hfill (\S2, \S7)
\item The geometric datum is assembled from a degenerate model; the lattice conditions on it hold by \lean{rfl} because the seed is defined as a copy of the row-side data; the map from the lattice to sections is zero. \hfill (\S3--\S5)
\item The lemmas that would extract geometric content are never called; one geometric field is vacuously true for the model; the degree equation is never proved. \hfill (\S6)
\item The height of the divisor datum is zero when the datum is built and is then overwritten by an explicit formula \lean{hExpl6}, so ``height $=$ formula'' holds by \lean{rfl};
two inequalities about that height are then proved by \lean{divClausesR_of} from about forty matrix estimates. \hfill (\S7)
\end{enumerate}
\textbf{Interpretation and requirements}
\begin{enumerate}[label=(B\arabic*),leftmargin=3.2em]
\item Hence the divisor side and the row side are one object, and the inequality that should trap a geometric height between a pairing and a main term compares a quantity with itself and holds by reflexivity. \hfill (Abstract, \S8)
\item Hence what the kernel certifies is only the tautology ``RiemannHypothesis $=$ RiemannHypothesis'' under the supplied definitions. \hfill (Abstract, \S8)
\item The argument that the source \emph{names} needs an independent geometric height. A formula that was written in to fit is not one, and a geometric argument cannot be both essential and trivially discharged. \hfill (\S8, \S9)
\item The algebraic zero-free step is downstream of the bundle. Remove the bundle and the proof term has no route to the conclusion; so the objection that a kernel-checked closed term is a proof does not apply. \hfill (\S10, App.~A)
\end{enumerate}
\textbf{Conclusion}
\begin{enumerate}[label=(C),leftmargin=3.2em]
\item The release does not contain a proof of the Riemann Hypothesis. \hfill (Abstract, \S9, \S13)
\end{enumerate}
\end{notebox}

Two features of this argument should be noted at once.
First, (A1) is conceded by \N: the dispute is not about axioms, \lean{sorry} or the identity of the statement.
Second, \N\ anticipates the obvious reply, namely that a kernel-checked closed term is a derivation of its statement whatever its intermediate steps look like, and answers it with (B4).
The weight of the paper therefore rests on (B1)--(B4). Section~\ref{sec:facts} checks (A1)--(A6); Section~\ref{sec:inference} examines (B1)--(B4).

\section{Q1: the observations are largely correct}
\label{sec:facts}

This section checks (A1)--(A6). The short answer is that \N\ read the code carefully and reports it accurately, and that the kernel-level export confirms the points on which a text search could mislead.
Table~\ref{tab:facts} at the end of the section collects the individual checks.
Three kinds of evidence are distinguished throughout: \textsf{S}, read by us in the source; \textsf{K}, derived from the kernel export through the dependency graph;
\textsf{G}, established by one of the independent Gemini agents (the letter names the agent; see Section~\ref{sec:agents}).

\subsection{The predicate and its wrapper (A3)}

\excerpt{anygeo}

Two predicates of the same type stand side by side. \lean{geoOf} asks for a positive integer $c$ with \lean{IsGeoOf} and for \lean{HasTrueSurf}; \lean{anyGeometry} returns \lean{True} on every input.
The predicate enters the development through a certificate \lean{OAt} with a field \lean{geo}, and through an inductive closure:

\excerpt{geocl}

The place where the \lean{geo} field is filled with \lean{trivial} is the one \N\ names:

\excerpt{divexprov}

\subsection{The route, and what lies on it (A2, A3)}
\label{sec:route}

\excerpt{main}
\excerpt{zfa}

The lines above are the last two links of the route described in \N, \S7. One level up, the witness \lean{h} is produced by \lean{nonempty_hyps_definite_side_of_closed_bound},
whose proof passes \texttt{geoCl anyGeometry} to three lemmas:

\excerpt{siegel_a}

A text search shows that an identifier is \emph{mentioned}; whether the proof term \emph{uses} it has to be read off the kernel export.

\begin{verdictbox}
\textbf{Kernel-level result (K).} In the dependency graph of the export, exactly \textbf{\NanyGeo} declarations depend on \lean{GlobalTheory.Mem.Top.Bl.anyGeometry}, and all five headline theorems are among them.
One chain is
{\small\lean{riemannHypothesis} $\to$ \lean{rhOfAll} $\to$ \lean{allChi} $\to$ \lean{UpperTheory.FinalTheorems.everyCharacter} $\to$ \lean{…everyCharacter_of_row}
$\to$ \lean{…every_character_of_literature_of_vWin} $\to$ \lean{UpperTheory.Lit.every_character_of_theta_integrals_local_matching}
$\to$ \lean{UpperTheory.Lit.nonempty_hyps_definite_side_of_closed_bound} $\to$ \lean{anyGeometry}.}

The names \lean{plugAll}, \lean{shapeWin}, \leanu{everyE₂}, \lean{linkDIVt_closedK} and \lean{geoOf} do not occur in the export: they are not used by any of the five theorems.
\end{verdictbox}

So (A2) and (A3) hold, including the correction that \N\ itself makes to earlier accounts: the route does not pass through \lean{plugAll} but through its twin \lean{zero_free_all_of_green_v_win}.

One detail of (A3) deserves a remark because \N\ leans on it (\S7, \S10, App.~A). The two arguments \lean{trivial}\ \lean{trivial} in \lean{zero_free_all_of_green_v_win} are not geometric data that has gone missing.
They instantiate two \emph{propositional parameters} of a generic lemma:

\excerpt{shapeof}
\excerpt{topshape}

\lean{forall_goal_of_shape_of_hyps} is stated for arbitrary propositions \lean{V} and \lean{NR}. The shape theorem used on the route is stated with \texttt{V := True} and \texttt{NR := True},
and its proof introduces the two corresponding hypotheses (\lean{hVC} and \lean{nR} in the source) and never refers to them again (\textsf{S}; confirmed independently by agent~H).
The route supplies \lean{trivial} because the statement being applied asks for a proof of \lean{True}.

\subsection{The degenerate model (A4)}
\label{sec:degenerate}

\excerpt{geodegen}

This is an instance of the abstract interface \texttt{DV.Geo F}: the variety is $\mathbb{N}$, and effectivity, nefness, bigness and the like are the proposition \lean{True}.
The file is called \path{Inhabited.lean}; \lean{mdlJunk}, \lean{cxJunk}, \lean{algDegenerate} and \lean{padDegenerate} are defined next to it.
Agent~E checked every row of the table in \N, \S3 (the fields of \lean{geoDegenerate}, of the datum \lean{divXFl} built on it, and of the seed \lean{x0}) against the source and found all of them accurate (\textsf{G-E}).
It also found that the description of \lean{graft} in \N\ is slightly simplified: besides the geometric fields from \lean{divXFl} and the lattice fields from the seed, about fifteen local fields come from a third datum.

\begin{verdictbox}
\textbf{Kernel-level result (K).} All of these objects are in the dependency cone of \lean{riemannHypothesis}. The numbers of declarations that depend on each are:
\lean{geoDegenerate} \NgeoDeg, \lean{algDegenerate} 41, \lean{padDegenerate} 42, \lean{mdlJunk} 44, \lean{cxJunk} 42, \lean{divXFl} 27, \lean{geoOver_stages} 25, the seed \lean{x0} 51, \lean{graft} 28.
The graph shows two further trivial instances of the interface that \N\ does not mention: an empty model \lean{geoExample} (53 dependents) and a one-point model \lean{pointGeo} (47 dependents).
The union of all these sets, together with the dependents of \lean{anyGeometry}, has \textbf{\Nflagged} elements.

The \emph{interface type} \lean{DV.Geo} is a different matter. It is a field of the record \lean{DivX} that is carried through the family data, and \NifaceDeps\ declarations depend on it.
But only \NaccDirect\ declarations refer directly to one of its field accessors, and only \NaccThmProof\ of these are theorems whose proof does so. No declaration eliminates a \lean{DV.Geo} or a \lean{DivX} record by its recursor.
\end{verdictbox}

The picture is therefore the one \N\ describes: a thin layer of declarations is written against an abstract geometry interface, and the release instantiates the interface with trivial models.
Section~\ref{sec:inference} asks what follows from that.

\subsection{The lattice side (A4)}

\excerpt{existsdivt}

The three lattice conditions of the existence statement \lean{DivEx} are discharged as \N, \S4 says: by \texttt{exists\_mjHTie … rfl HEq.rfl}, by \texttt{fun \_ => rfl}, and by \texttt{lamAre\_of\_eq … rfl} (\textsf{S}, \textsf{G-E}).
They can be, because the seed sets \texttt{MjH := mjH Z}, \texttt{logLam := logLamOf Z hH} and \texttt{kapR := fun \_ => kapSharp D}.

Agent~E added a point that \N\ passes over quickly and that matters for (B1). The objects that are copied are not themselves trivial.
\lean{mjLat}, \lean{logLamOf} and \lean{kapSharp} are Hermitian lattices, logarithms of eigenvalues and a spectral constant computed from the arithmetic (``row-side'') data,
and the release proves lemmas with content about them (\lean{mjH_good}, \lean{jM_lat}, \lean{jM_metric}) from hypotheses \lean{Holds} and \lean{Pins} (\textsf{G-E}; all six names are in the export, \textsf{K}).
What holds by \lean{rfl} is that the datum \emph{carries} these lattices, not anything about the lattices.

\subsection{The height (A6)}
\label{sec:height-facts}

\excerpt{hexpl6}
\excerpt{heighteq}

The height slot of the datum is overwritten by \lean{setDiv} with \texttt{p.ht}, which is the formula \lean{hExpl6}; the field \lean{h_eq} then holds by \lean{rfl}. This is (A6), and it is accurate (\textsf{S}, \textsf{G-G}).
Agent~G checked the remaining details of \N, \S7 and reports the following (we re-checked the items marked \textsf{S}):
\begin{itemize}
\item Every caller of \lean{exists_divStage_of_pk} passes the zero function as initial height, as \N\ says (\textsf{G-G}).
\item Of the four sites at which \N\ says \lean{h_eq} is constructed, two (\path{DisjointSplitLevelPlaces.lean}, lines 331 and 501) concern a field of the same name in a different structure,
\texttt{h\_eq : h = hP * k} at line 264 of that file (\textsf{S}); two other sites that propagate the real field are not listed in \N.
\item The construction \lean{linkDIVt_closedK}, which \N\ uses to explain how the height is supplied, is not in the export (\textsf{K}). \N\ notes later in \S7 that the live route uses \lean{closed_linkDIVt_of_closed_bridge_kernel} instead; that theorem is in the export.
\item On the live route the pairing \lean{Bk} is not identically zero: it is tied to the period pairing \lean{Bu} by a proved rationality statement \lean{BuRational}, and the case $\mtt{Bk}\ u \neq 0$ is used downstream (\textsf{G-G}; the lemmas named by the agent are in the export, \textsf{K}).
\end{itemize}
The last sentence of (A6), that the two inequalities about the height are proved from matrix estimates, is also accurate. It is examined in Section~\ref{sec:squeeze} because it decides (B1).

\subsection{The dead extraction lemmas (A5)}
\label{sec:dead}

\N, \S6 reports that \lean{GeoCl} is never opened by case analysis or induction and that six lemmas written to extract geometric content occur only at their own declarations.
A text search cannot quite establish the first claim, since pattern matching need not name the type. The kernel export can.

\begin{verdictbox}
\textbf{Kernel-level result (K).}
\begin{itemize}
\item The recursor \lean{GeoCl.rec} is referenced by no declaration outside the inductive block of \lean{GeoCl} itself; \lean{GeoCl.casesOn} and \lean{GeoCl.recOn} do not occur in the export.
Any case analysis or induction on a \lean{GeoCl} proof compiles to the recursor. \emph{No proof in the cone ever opens a} \lean{GeoCl} \emph{certificate.}
\item The constructor \lean{GeoCl.base} is used once (in \lean{geoCl_of}) and \lean{GeoCl.step} once (in \lean{oAt_setDiv}).
\item The \lean{geo} field of the certificate \lean{OAt} is read by exactly two theorems, \lean{divEx_geoCl} and \lean{oAt_setDiv}, and each passes it straight into a \lean{GeoCl} constructor. \lean{OAt.rec} has no referrer.
By contrast the other field, \lean{lam}, is also read by \lean{GlobalTheory.T.parameter_of_discrete_series_ray_alt}: in the same certificate one field has a consumer and the other has none.
\item \lean{geoOf}, \lean{GeoTiedP.geoTied}, \lean{GeoTiedP.formsTieX}, \lean{degree_value_of_determinant} and \lean{SlopeBoundFamilyStatement} do not occur in the export.
\end{itemize}
\end{verdictbox}

The last item bears on the remark in \N, \S6 that the degree equation is ``never proved''. It occurs only as a hypothesis inside \lean{SlopeBoundFamilyStatement} and one admissibility bundle.
Since \lean{SlopeBoundFamilyStatement} is not in the export, no declaration used by the final theorem assumes that equation. An unproved hypothesis of an unused statement is not a gap. \Fnote

\subsection{Summary of the fact check}

{\small
\renewcommand{\arraystretch}{1.18}
\begin{longtable}{@{}r>{\raggedright\arraybackslash}p{6.5cm}c>{\raggedright\arraybackslash}p{6.3cm}@{}}
\caption{Observations of \N\ and what we found. Evidence: \textsf{S} source, \textsf{K} kernel export, \textsf{G-x} Gemini agent x.}\label{tab:facts}\\
\toprule
\# & Observation in \N\ (section) & Finding & Evidence and remarks \\
\midrule
\endfirsthead
\toprule
\# & Observation in \N\ (section) & Finding & Evidence and remarks \\
\midrule
\endhead
\bottomrule
\endfoot
1 & Closed theorem of Mathlib's \lean{RiemannHypothesis}; no axiom, \lean{sorry}, \lean{native_decide}, \lean{unsafe}, \lean{opaque} (Abstract, \S12) & \OK & \textsf{S}, \textsf{K}: three standard axioms only; statement cone is Mathlib only (Section~\ref{sec:taut}) \\
2 & Route to the final theorem; \lean{plugAll} not on it (\S1, \S7, \S11) & \OK & \textsf{S}, \textsf{K} \\
3 & \texttt{anyGeometry := fun \_ \_ \_ => True}; 118 uses in 24 files (\S2) & \OK & \textsf{S}; \textsf{K}: \NanyGeo\ dependents, final theorem included \\
4 & \lean{divEx_prov} fills the geometric field with \lean{trivial} (\S2) & \OK & \textsf{S}, \textsf{G-E} \\
5 & \lean{divEx_of_ox} ignores three hypotheses named after published results (\S2) & \OK & \textsf{S}: body is \texttt{fun \_ \_ \_ 𝒮 hOx => divEx\_prov 𝒮 hOx}. \textsf{G-E}: at the call sites the three arguments are theorems proved in the release \\
6 & Field table of the degenerate model, of \lean{divXFl} and of the seed (\S3) & \OK & \textsf{G-E}: every row verified; \textsf{K}: all in the cone \\
7 & \lean{graft} takes geometric fields from \lean{divXFl}, lattice fields from the seed (\S3) & \PARTLY & \textsf{G-E}: also about 15 local fields from a third datum \\
8 & Lattice conditions of \lean{DivEx} discharged by \lean{rfl} (\S4) & \OK & \textsf{S}, \textsf{G-E}; the copied lattices are substantive objects with proved lemmas \\
9 & Row-side conditions \lean{Holds}, \lean{InGeo} have content (\S5) & \OK & \N's own statement; \textsf{K}: in the cone\Frow \\
10 & \lean{GeoTiedP.geoTied}, \lean{.formsTieX} never called (\S5) & \OK & \textsf{K}: neither name occurs in the export \\
11 & \lean{GeoCl} never opened (\S6) & \OK & \textsf{K}: \lean{GeoCl.rec} has no referrer outside its own block \\
12 & Six extraction lemmas occur once (\S6) & \OK & \textsf{K}: those checked are absent from the export\Fsix \\
13 & Degree equation only a hypothesis, never proved (\S6) & \OK & \textsf{K}: and the statement containing it is not used by the final theorem \\
14 & Two shape inputs supplied by \lean{trivial} (\S7) & \OK & \textsf{S}, \textsf{G-H}: they instantiate two unused propositional parameters at \lean{True} \\
15 & Height overwritten by \lean{hExpl6}; \lean{h_eq} by \lean{rfl} (\S7) & \OK & \textsf{S}, \textsf{G-G} \\
16 & \lean{h_eq} constructed at four listed sites (\S7) & \PARTLY & \textsf{S}, \textsf{G-G}: two of the four are a different field of the same name \\
17 & Initial height is zero (\S7) & \OK & \textsf{G-G} \\
18 & \lean{WordFacts}: two inequalities, proved by \lean{divClausesR_of} from about forty inputs (\S7) & \OK & \textsf{S}, \textsf{K}, \textsf{G-G}: \NdivInputs\ fields; \lean{WordFacts} has a third conjunct not shown in \N \\
19 & \lean{Bk} is the pairing or else 0 (\S7) & \PARTLY & \textsf{G-G}: on the live route it is tied to the pairing by a proved statement \\
20 & The squeeze ``compares a quantity with itself and holds by reflexivity'' (Abstract, \S8) & \NO & \textsf{S}, \textsf{K}, \textsf{G-G}, \textsf{G-H}: contradicted by row 18; Section~\ref{sec:squeeze} \\
21 & Only the tautology ``RiemannHypothesis $=$ RiemannHypothesis'' is certified (Abstract, \S8) & \NO & \textsf{K}, \textsf{G-H}: Section~\ref{sec:taut} \\
\end{longtable}
}

\section{Q2: the conclusion does not follow}
\label{sec:inference}

\subsection{What a kernel certificate says}
\label{sec:kernel}

Three elementary facts about Lean's logic carry most of the argument. None is specific to this release.

\begin{proposition}[Definitions are conservative]\label{prop:cons}
A definition \texttt{def c : T := v} names a closed term $v$ that has already been type-checked. If a proof term $t$ of a statement $P$ mentions $c$ and $P$ does not, then replacing $c$ by $v$ throughout $t$ gives a proof of $P$ that does not mention $c$.
Hence no choice of auxiliary definitions can make provable a statement that was not provable without them.
\end{proposition}

\begin{proposition}[Vacuity needs a hypothesis]\label{prop:vac}
A statement is vacuously true when it has the form $H \to P$ and $H$ cannot be met, or when $H$ has been weakened until it no longer bears on $P$. A proof of such a statement says nothing about $P$.
A statement with no hypothesis leaves no room for this: a proof of $P$ is a proof of $P$.
Intermediate lemmas do have hypotheses, but in a closed proof every hypothesis of every lemma that is applied is discharged by a checked proof term.
That a hypothesis was discharged \emph{easily} does not mean it was not discharged.
\end{proposition}

\begin{proposition}[The four exits]\label{prop:exits}
Suppose a checker accepts a closed term $t$ with $t : \mtt{RiemannHypothesis}$ whose dependencies contain only the axioms \lean{propext}, \lean{Quot.sound} and \lean{Classical.choice}.
Then the Riemann Hypothesis, as formalized in Mathlib, has been derived, unless one of the following holds:
\begin{enumerate}[label=(\roman*)]
\item the derivation uses a further axiom or \lean{sorry} after all;
\item the statement proved is not the Riemann Hypothesis (a hijacked name, a modified library, a wrong definition in Mathlib);
\item the checker is unsound and the term exploits this;
\item the object that was checked is not the object that is claimed.
\end{enumerate}
\end{proposition}

The argument of \N\ takes none of these exits. It grants that the term type-checks against Mathlib's statement without extra axioms (A1) and then argues from the \emph{content of intermediate definitions}.
By Proposition~\ref{prop:cons} that content cannot matter to whether the final statement has been derived.
\N\ is aware of this reply and rejects it in Appendix~A; we come to that in Section~\ref{sec:remove}. First we examine the two descriptions on which the title claim of \N\ rests.

\subsection{``Only a tautology'' (B2)}
\label{sec:taut}

\N\ states that the formalization certifies only the tautology ``RiemannHypothesis $=$ RiemannHypothesis'' under the supplied definitions, and that the kernel accepts the term because a term that is definitionally an identity always type-checks (Abstract, \S8).
As a description of the checked artefact this is wrong on each count.

\begin{itemize}
\item \textbf{The type is not an equation.} The certified declaration is \texttt{theorem riemannHypothesis : RiemannHypothesis}. Its type is the proposition \lean{RiemannHypothesis}, which unfolds in Mathlib to the statement that every zero of \lean{riemannZeta} other than the trivial ones and $s = 1$ has real part $1/2$.
The proposition $\mtt{RiemannHypothesis} = \mtt{RiemannHypothesis}$ is a different proposition, proved by \lean{rfl}; no term of that kind has type \lean{RiemannHypothesis}.
\item \textbf{The statement does not involve the supplied definitions.} In the dependency graph, the statement cone of \lean{riemannHypothesis} (Definition~\ref{def:cones}) consists of \NstmtMain\ declarations.
All of them are library declarations; none was declared by the release (\textsf{K}). The phrase ``under the supplied definitions'' would be apt for a theorem whose \emph{statement} mentioned \lean{anyGeometry} or \lean{hExpl6}. This one mentions neither.
\item \textbf{The statement is the one in Mathlib.} We rebuilt the export of the five statements from official Lean 4.33.1, the official Mathlib cache and nothing else, and obtained a file identical byte for byte to the published challenge export (Section~\ref{sec:rebuild}).
The checker \texttt{con-leche}, in comparison mode, accepts the full proof export as a justification of exactly those statements (Section~\ref{sec:conleche}).
\end{itemize}

Whatever one thinks of the route, the destination is fixed, and it is not an identity.

\subsection{``Compares a quantity with itself'' (B1)}
\label{sec:squeeze}

The central factual claim of \N\ is that, because of the three definitional steps, the inequality at the heart of the argument compares a quantity with itself and holds by reflexivity (Abstract, \S8, \S10).
The statement in question is the following; it has three conjuncts, of which \N\ quotes the second and third.

\excerpt{wordfacts}

The second conjunct bounds the height of the pairing \lean{Bk} from above by \texttt{D.div.h u}. The third bounds $\mu^\vee + h - F_o$ from above by the main term plus error terms.
After substituting the installed value $h = \mtt{hExpl6}$ they read
\[
\mathrm{ht}\bigl(\mtt{Bk}\,u\bigr) \;\le\; \mtt{hExpl6}(u)
\qquad\text{and}\qquad
\mu^\vee + \mtt{hExpl6}(u) - F_o \;\le\; \mathrm{MAIN}(u) + \text{error}.
\]
Neither has the form $X \le X$: the left side of the first is a mean of logarithms of operator norms of the pairing, the right side of the second is the main term, and \lean{hExpl6} is a sum of explicit size and weight terms (Section~\ref{sec:height-facts}).
Nor are they proved by reflexivity. \lean{WordFacts} is established by \lean{wordFacts_of} from separately proved clauses:

\excerpt{wordfactsof}

The hypothesis \lean{tH} is where ``height $=$ formula'' enters; it is the only ingredient that comes from \lean{rfl}. The clauses \lean{c2} and \lean{c3} are the two inequalities, and they come from here:

\excerpt{divclauses}
\excerpt{divclausesof}

The proof of the third clause occupies lines 210--245 of the same file; it invokes, among others,

\excerpt{divclausesof2}
\excerpt{divclausesof3}

The input record \lean{DivInputsR} has \NdivInputs\ fields: lattice and frame conditions, bounds on operator norms along a chain of matrices, an eigenvalue bound, valuation bounds at finitely many places, a covolume-type bound and numerical side conditions (\textsf{S}; listed field by field by agent~G).
Agent~G traced where such a record is constructed on the live route and examined thirteen of its fields; it reports that each is supplied by a lemma with a substantive proof and none by \lean{rfl}, \lean{trivial} or elimination of an empty type (\textsf{G-G}).

\begin{verdictbox}
\textbf{Kernel-level result (K).} The dependency cone of \lean{Formalism.L2.DIV.divClausesR_of} contains \NdivClausesAll\ declarations, of which \NdivClausesProj\ belong to the release and \textbf{\NdivClausesThm} of those are theorems.
It contains none of the objects flagged in \N: not \lean{anyGeometry}, not the degenerate or junk models, not \lean{divXFl}. A proof by reflexivity has an empty cone.
\end{verdictbox}

This is also what \N\ itself says in \S7, where the two inequalities are attributed to \lean{divClausesR_of} and described as estimates on matrices that make no use of the geometric certificate, and in the table of \S8, where ``height $=$ formula'' is marked as true by \lean{rfl} but the two inequalities are not.
The Abstract and \S8 of \N\ thus assert something that \S7 of \N\ and the code both contradict.
What holds by \lean{rfl} is that a slot contains a formula. The inequalities about that formula are theorems.

\subsection{``Assigned, not derived'' (B3)}
\label{sec:assigned}

Once (B1) is set aside, what \N\ has shown about the height is this: the quantity $h$ that stands between the pairing and the main term is not introduced as the height of a geometric object and then computed; it is introduced as a formula, and the two bounds are proved for the formula.
\N\ objects that the height was chosen to fit, and that inequalities about a chosen formula say nothing about a geometric height (\S8).

The second half of that sentence is true, and irrelevant to validity. To prove $A \le C$ it suffices to exhibit some $h$ with $A \le h$ and $h \le C$.
The logic asks nothing about where $h$ came from; in particular it does not ask that $h$ be ``independent'' of $A$ and $C$ or that it have a geometric meaning.
Choosing the intermediate quantity so that both bounds can be proved is what proving such a statement consists of, and the bounds still have to be proved.
They are not free: with a different value in the slot, for instance the initial value $0$, the same two lines would be different statements, and the release does not prove those.

\N\ concedes the essential point in \S8 when it observes that adding the two inequalities cancels $h$ and leaves a comparison between the pairing and the main term, established by \lean{divClausesR_of}.
That comparison is a statement about the arithmetic data alone, it is what the downstream argument consumes, and it has been proved.
The remaining objection is that this is not the argument that the identifiers and the directory names of the release announce.
That is so (Section~\ref{sec:credit}). But a derivation is not invalidated by the names of its lemmas.

The accompanying file \path{evidence/toy/ToyInterpolant.lean} reproduces the pattern in a dozen lines of Lean: a record with a height slot, an overwrite, \texttt{h\_eq := rfl}, two honest bounds for the installed formula,
and a closed conclusion in which the height no longer appears. It also shows the bounds failing for the slot's initial value.

Finally, \N, \S9 says that a geometric argument cannot be both essential and trivially discharged. We agree with the sentence. Together with the fact that the geometric slot \emph{is} trivially discharged in a proof the kernel accepts, it yields by contraposition that the geometric slot is not essential.
\N\ draws the opposite conclusion by taking essentiality as a premise, which is the point at issue.

\subsection{``Remove the package and no route remains'' (B4)}
\label{sec:remove}

Appendix~A of \N\ addresses the reply based on Proposition~\ref{prop:exits}. It says that the reply assumes that removing the vacuous inputs leaves an independent derivation behind,
and that this is not so, since the algebraic step consumes the package assembled from \texttt{geoCl anyGeometry} and the two \lean{trivial} inputs (also \S10).

The reply assumes no such thing. It says that the term as it stands, with those inputs in place, is a derivation. Nothing has to be removed, and no second route is needed.

The test proposed in \N, deleting an input and observing that the term no longer reaches its conclusion, does not measure soundness.
Every argument of every lemma application fails it, because a term with a missing argument is not a term. Agent~H gave a minimal example:
\begin{Verbatim}[fontsize=\small,frame=single,rulecolor=\color{gray!50}]
structure Wrap (P : Prop) where
  dummy : P
  real  : 2 + 2 = 4
theorem t : 2 + 2 = 4 := (Wrap.mk (P := True) trivial rfl).real
\end{Verbatim}
\noindent Delete the package \texttt{Wrap.mk …} and \texttt{t} has no route to its conclusion. Yet \texttt{t} is sound and does not depend, in any sense that matters, on \texttt{dummy : True}.

The test that does measure whether a derivation relies on the \emph{content} of an input is replacement: would the derivation go through with a different input of the same type?
For the inputs that \N\ calls vacuous the answer can be read off the statements of the consumers.

\paragraph{The geometry predicate.} The lemmas on the route that receive \texttt{geoCl anyGeometry} are stated for an arbitrary predicate $\Gamma$. For example:

\excerpt{mono_sig}

Agent~H lists further consumers on the route of the same form, among them \lean{rowFun_at_gridA_alt}, \lean{linked_bundles_l3_of_open_alt} and \lean{allRows_of_open0_alt} (\textsf{G-H}; we confirmed the quantifier $\forall\,\Gamma$ in the three statements, \textsf{S}, and their presence in the export, \textsf{K}); agent~B of the first round reached the same conclusion for the whole route.
A lemma proved for all $\Gamma$ cannot use any property of $\Gamma$. Applying it to $\Gamma = \mtt{True}$ is as legitimate as applying it to any other value, and it would be equally provable for any other value.

\paragraph{The two shape inputs.} As shown in Section~\ref{sec:route}, they are proofs of \lean{True} demanded by a lemma that is generic in two propositions \lean{V} and \lean{NR} and by a shape theorem stated for $\mtt{V} = \mtt{NR} = \mtt{True}$.

\paragraph{The bundle.} The object handed to the algebraic step is a proof of the proposition \texttt{Nonempty (HypsAlt …)}.
The type \lean{HypsAlt} does not mention the geometry predicate: its dependency cone contains neither \lean{GeoP} nor \lean{GeoCl} nor \lean{anyGeometry} (\textsf{K}); it does contain the interface type \lean{DV.Geo}, through the family records, but none of its trivial instances.
Since \texttt{Nonempty (…)} is a proposition, proofs of it are indistinguishable to anything that consumes them. The algebraic step cannot depend on \emph{how} the bundle was shown to be inhabited, only on the fact that it is.

\excerpt{hypsalt}

\paragraph{Upstream and downstream.} \N\ describes the algebraic estimates as downstream of the vacuous package. In the order of application that is right. In the order of dependency it is the reverse, and the dependency order is the one that matters for what has been proved.

\begin{verdictbox}
\textbf{Kernel-level result (K).}
\begin{itemize}
\item \lean{UpperTheory.Top.every_character_of_green_v_win}, which derives the conclusion about all Dirichlet characters from \emph{any} proof of \texttt{Nonempty (HypsAlt …)}, has a dependency cone of \NgreenAll\ declarations,
\NgreenProj\ of them from the release (\NgreenThm\ theorems). The cone contains none of the flagged objects. The algebraic step was proved without them.
\item The theorem that supplies the bundle, \lean{nonempty_hyps_definite_side_of_closed_bound}, has a cone of \NpackageAll\ declarations, \NpackageProj\ from the release, of which \textbf{\NpackageThm} are theorems.
Of all \NconeProj\ release declarations in the cone of \lean{riemannHypothesis}, \Nflagged\ depend on any flagged object and \textbf{\NfreeProj} do not.
\end{itemize}
\end{verdictbox}

The package is thus not vacuous. Inhabiting \lean{HypsAlt} takes more than six hundred thousand theorems; a predicate set to \lean{True} occupies one parameter slot along the way, in lemmas that hold for every value of that slot.

\subsection{Why the intuition misleads}
\label{sec:asym}

The feeling that setting a condition to \lean{True} is a way of cheating is sound when only one party is checked. In a proof a predicate always faces two.
The \emph{producer}, who must show that the predicate holds, gains when it is weakened; at \lean{True} the producer writes \lean{trivial}.
The \emph{consumer}, who must derive something from it, loses exactly as much; from \lean{True} the consumer learns nothing.
For weakening to pay, the producer must gain while the consumer does not lose. In Lean that happens in three ways: the consumer does not exist (the weakened predicate is left in the \emph{final} statement);
the consumer is not checked (an axiom or \lean{sorry}); or the final statement is exchanged for another. With the final statement fixed and no axioms, the consumer has to reach the conclusion without the content of the predicate, and then the predicate was not needed.

The file \path{evidence/toy/ToyGeometry.lean} shows the three situations side by side: a closed theorem obtained through a predicate set to \lean{True}, with no axioms;
a consumer that really needs the predicate, for which setting it to \lean{True} makes the required lemma \emph{false}; and the cheating variants, each caught by \texttt{\#print axioms} or visible in the statement.
The objection of \N\ would be correct against the third kind of development. The release is of the first kind.

\section{An independent search for real defects}
\label{sec:defects}

Since the argument of \N\ cannot establish invalidity, one may still ask whether the release is invalid for some other reason.
Proposition~\ref{prop:exits} lists the possibilities, and each can be checked mechanically. This section does not depend on \N.

\subsection{Exit (i): axioms, \texttt{sorry}, trusting the compiler}

A single-pass scan of all 175,187 source files gives: 0 \lean{axiom} declarations; 0 occurrences of \lean{native_decide}, \lean{ofReduceBool}, \lean{implemented_by}, \lean{extern}, \lean{unsafe}, \lean{opaque}, \lean{csimp};
5 occurrences of \lean{sorry}, all in \path{src/Challenge.lean}, a file added by the community for the comparison tool, absent from the author's archive and not imported by \path{Main.lean}.
The kernel export gives the same answer without relying on text: the only axiom records are \lean{propext}, \lean{Quot.sound} and \lean{Classical.choice}; no record is flagged unsafe or partial;
the names \lean{sorryAx}, \lean{Lean.ofReduceBool}, \lean{Lean.ofReduceNat} and \lean{Lean.trustCompiler} do not occur; all 91 opaque constants are library declarations.
This agrees with \N\ (\S12) and with the first-round hygiene scan of agent~D.

\subsection{Exit (ii): a substituted statement}
\label{sec:rebuild}

The five theorems in \path{Main.lean} have no hypotheses. No declaration of the release is named \lean{RiemannHypothesis}, \lean{riemannZeta} or \lean{DirichletCharacter}.
All 1,695 \lean{notation} commands in the release are \lean{local} and one is \lean{scoped}; there is no global notation and no global macro other than tactic macros, so the parsing of \path{Main.lean} cannot have been altered.
Mathlib is pinned at the official tag \texttt{v4.33.1} with no path override.

The strongest check is a reconstruction. The comparison of a proof export with a statement export is only as good as the statement export, so we rebuilt it:
official Lean 4.33.1 binaries; a project containing only the five statements with \lean{sorry} and \texttt{import Mathlib}; 8,690 precompiled files from the official Mathlib cache; \texttt{lean4export} at tag \texttt{v4.33.0}.

\begin{verdictbox}
The rebuilt file has 350,862,311 bytes and SHA-256 \texttt{bd484ac9…bdd6ab2b9}. It is \textbf{byte-identical} to the published \path{challenge.ndjson}.
The reference against which the proof is compared is therefore determined by official Mathlib and the five statements alone.
\end{verdictbox}

\subsection{Exit (iii): the kernel}
\label{sec:conleche}

The release contains metaprogramming that deserves a look, and this is where a real exploit would have to hide.
The files \path{BuildAid/D1..D5.lean} call \lean{Lean.modifyEnv}, but only to extend a deny-list for library suggestions.
A command \lean{cpt_copy_thms} copies declarations in bulk through the standard \lean{Lean.addDecl}, so each copy passes the kernel.
Two custom tactics, \lean{kexactB} (234 uses) and \lean{kcopy}, assign a term to the goal without elaborator-side type checking; \lean{kcopy} builds raw projection terms:

\excerpt{kcopy}

This is legitimate engineering, but it means that those proofs are guarded by the kernel alone. If a kernel had a defect in the checking of projections, this is where it could be exploited.
Reading the source cannot settle that; an independent replay can.

We built \texttt{con-leche}~\cite{conleche}, a checker with its own term representation and a formalized consistency argument, from source at commit \texttt{bc667610} and ran
\begin{Verbatim}[fontsize=\small,frame=single,rulecolor=\color{gray!50}]
con-leche --verified --compare challenge.ndjson solution.ndjson
\end{Verbatim}
\noindent on the complete 61.57\,GB export. The run ended with exit code 0 and this output:
\VerbatimInput[fontsize=\scriptsize,frame=single,rulecolor=\color{gray!50},breaklines=true,breakanywhere=true]{../evidence/con-leche/compare.stdout.log}
\noindent All 1,468,418 declarations of the proof export were accepted, and each of the five statements of the challenge export was found to be justified by a constant of the same name and type whose statement reads the same declarations on both sides.
The run took 73 minutes on 16 cores with a peak of 47.4\,GB of memory. The verification page~\cite{mit} reports the same counts for \texttt{con-leche} and passes by further independent kernels, which we did not re-run.

\subsection{Exit (iv): source versus checked artefact}

The hashes in Table~\ref{tab:prov} tie the archive, the mirror and the two exports to the published values.
The one link we did not reproduce is the build that turns the 175,187 source files into the proof export; for it we rely on the sandboxed rebuild reported on the verification page.
This limitation affects the claim that the \emph{source} proves the theorem. It does not affect the claim that the \emph{export} does: the export was checked directly, and the dependency analysis of this note was performed on it.

\subsection{Result}

None of the four exits shows any sign of being open.

\section{Independent verification by Gemini agents}
\label{sec:agents}

A substantial part of the checking in this note was not done by the main analysis but by a second, independent line of verification, and the note would be weaker without it.
Eight agents were run through the \texttt{agy} command-line tool with Google's Gemini~4 Neon model at its highest reasoning setting.
This is a different model family from the system used for the main analysis (Section~\ref{sec:contrib}), which is the reason for using it: an error of reading or of reasoning made by one is less likely to be repeated by the other.

Each agent started from an empty context. It received one sub-question, read-only access to the source tree and, in the second round, the text of \N.
It did not see our notes, the dependency graph, or any conclusion of ours, and it was instructed to cite file and line for every claim, to answer ``not verified'' rather than guess, and to report what the code says whichever side that favours.
Four agents (A--D) worked on the claims of \N\ as publicly summarized; four (E--H) worked on the full text. The prompts and the complete reports are in the repository.
Six further agents (I--O) were run for the addendum and are described there (Section~\ref{sec:agents3}).

\begin{table}[ht]
\centering\small
\renewcommand{\arraystretch}{1.22}
\begin{tabularx}{\textwidth}{@{}c>{\raggedright\arraybackslash}p{3.0cm}>{\raggedright\arraybackslash}X@{}}
\toprule
Agent & Sub-question & What it established \\
\midrule
A & declaration-level dependency trace & Found by source reading the chain from \lean{riemannHypothesis} to \lean{anyGeometry}, and that \lean{plugAll}, \leanu{everyE₂} and \lean{geoOf} are off it. The kernel graph later gave the same answer. \\
B & logical role of the predicate $\Gamma$ & Every consumer on the route is stated for arbitrary $\Gamma$; no elimination of \lean{OAt.geo} or \lean{GeoCl}; the final bundle type does not contain $\Gamma$. \\
C & the secondary claims & \textbf{First located} \lean{geoDegenerate}, \lean{mdlJunk} and the related definitions, which a public review had searched for and not found. \\
D & hygiene scan & No axioms, no \lean{native_decide}, \lean{sorry} only in the community file; \textbf{first pointed out} that \lean{kexactB} and \lean{kcopy} make an independent kernel replay necessary. \\
\midrule
E & \N, \S2--\S4 & Verified every row of the model table; found that the three discarded hypotheses of \lean{divEx_of_ox} are theorems proved in the release; \textbf{showed that the ``copied'' lattices are substantive objects} with proved lemmas. \\
F & \N, \S5--\S6 & \Fsummary \\
G & \N, \S7--\S8 & Traced \lean{wordFacts_of} to \lean{divClausesR_of} and listed all \NdivInputs\ fields of the input record; audited thirteen of them on the live route; \textbf{found two mis-cited sites in \N}; \textbf{found that} \lean{Bk} \textbf{is tied to the period pairing by a proved statement}. Its answer to the key question: the inequalities are substantive, not reflexive. \\
H & the argument of \N\ (Abstract, \S8--\S10, App.~A) & An independent reconstruction of the argument in thirteen steps and an assessment of each; the \texttt{Wrap} example of Section~\ref{sec:remove}; a list of universally quantified consumers; the contraposition of Section~\ref{sec:assigned}; a list of what would constitute a real defect and of what legitimately follows from \N. \\
\bottomrule
\end{tabularx}
\caption{The eight independent agents (Gemini~4 Neon). Bold marks findings that originated with the agent and were then adopted here.}
\label{tab:agents}
\end{table}

All eight agents reached the same overall assessment as the main analysis: the observations of \N\ about the code are largely accurate, and the inference to invalidity does not hold.
Agreement between two systems is evidence, not proof, and three qualifications apply.
The agents and the main analysis read the same source tree, so this is a cross-check of reading and reasoning, not an independent source of the code.
The prompts were written by us; the prompt of agent~B contained a hint towards its answer, and agent~H was given as premises the four facts that \N\ itself concedes or that Section~\ref{sec:facts} establishes.
And an agent's statement is used in this note as a fact only where it was re-checked in the source or in the kernel graph; otherwise it is labelled \textsf{G} in Table~\ref{tab:facts}.

\section{What the paper under review does establish}
\label{sec:credit}

Rejecting the conclusion of \N\ should not obscure what its reading of the code has achieved. Four things follow from it, and from the kernel-level confirmation given here.

\begin{enumerate}
\item \textbf{The formal proof does not use the geometry that the release is named after.}
The predicate \lean{geoOf}, the conditions \lean{IsGeoOf} and \lean{HasTrueSurf} and the lemmas around them are in the source and outside the dependency cone of the final theorem.
The interface \lean{DV.Geo} is instantiated only by trivial models. Whatever the proof term proves, it does not prove it by way of heights on an arithmetic variety supplied through this interface.
\item \textbf{The height in the central inequality is an explicit formula.} The release contains no statement identifying \lean{hExpl6} with the height of a geometric object.
If such an identification belongs to the idea behind the proof, it is not part of what was formalized. Equally, the formal proof does not need it.
\item \textbf{Names are not evidence, in either direction.} A directory called \path{AbelianVarietyTheory} does not show that the proof uses abelian varieties; this is the valuable half of the lesson of \N.
An identifier called \lean{mdlJunk} does not show that the proof is junk; that is the other half. Only terms and types decide.
\item \textbf{An explanation is still owed.} A theorem of this importance whose proof is certified but not yet understood is in an unusual state, and the map drawn in \N\ of which parts of the release are live and which are dead is a real contribution to changing that.
As agent~H put it, the fitting response to these findings is not to declare the theorem unproved but to strip the unused scaffolding and to set out the argument that the proof term actually carries out.
\N\ cites an exposition of the proof by Jaswal~\cite{jaswal}; we have not examined it and take no position on it.
\end{enumerate}

These are statements about exposition and about the relation between a formal proof and an informal narrative. They are compatible with the theorem being proved, and on the evidence of Sections~\ref{sec:inference} and~\ref{sec:defects} it is.

\paragraph{A remark on ``long chains that pass Lean''.}
The episode prompted the informal claim that any language model can be made to output a sufficiently long chain of steps that passes Lean while only pretending to reason.
A chain that passes the kernel, ends in the intended statement and uses no extra axiom is a proof; the kernel does not grade sincerity.
What a model can fake is one of the four exits of Proposition~\ref{prop:exits}, and each of them is detectable by the mechanical checks of Section~\ref{sec:defects}.

\section{Limits of this audit}
\label{sec:limits}

\begin{enumerate}
\item \textbf{The mathematics was not read.} This note establishes that a particular objection fails and that standard mechanical checks pass. It does not explain the proof and is not a substitute for mathematical review.
\item \textbf{The trust base remains.} The conclusion rests on the consistency of Lean's type theory with the three standard axioms, on the correctness of at least one of the checkers that accepted the export
(we ran one, \texttt{con-leche}, ourselves; others are reported on the verification page), on the faithfulness of \texttt{lean4export}, and on the adequacy of Mathlib's definitions of the Riemann zeta function and of \lean{RiemannHypothesis}.
None of these is questioned in \N, and we did not audit them beyond what Section~\ref{sec:defects} describes.
\item \textbf{The source-to-export build was not reproduced.} See Section~\ref{sec:defects}, exit~(iv).
\item \textbf{A dependency graph is syntactic.} It tells exactly which declarations a term refers to. Absence from the export means ``not used''; presence means ``referred to'', which is weaker than ``essential''.
The split between library and release is by name against the Mathlib environment and may misplace a small number of automatically generated auxiliary declarations.
\item \textbf{Counts measure extent, not depth.} That the bundle witness depends on \NpackageThm\ theorems shows that it is not obtained from nothing. It does not show that those theorems are deep; we did not grade them.
\item \textbf{Not every agent finding was re-verified.} Those that were not are marked \textsf{G} without \textsf{S} or \textsf{K} in Table~\ref{tab:facts}.
\item \textbf{The version of \N.} We reviewed the 17-page version with the hash given in Table~\ref{tab:prov}. \N\ has been revised more than once, and a later version may differ.
\end{enumerate}

\section{Conclusion}
\label{sec:conclusion}

The paper under review reads the Lean release closely and reports, correctly, that its geometric layer is hollow: a predicate that is \lean{True}, an interface instantiated by trivial models, lattice conditions that hold because the datum was built as a copy, and a height that is a formula.
We have confirmed these observations in the source and, where it matters, in the kernel-level export of the proof.

The conclusion that the release contains no proof of the Riemann Hypothesis does not follow from them.
It rests on two descriptions of the code that are false (that the central inequality holds by reflexivity, and that only a tautology is certified) and on a test of soundness, deletion of an input, that every valid derivation fails.
The certified object is a closed term of Mathlib's statement with the three standard axioms. The inequalities are theorems with \NdivClausesThm\ theorems behind them. The consumers of the trivial inputs are quantified over all inputs.
The hollow layer is a part of the proof term that the proof does not need.

The title of this note is deliberately symmetrical with the title of the paper it answers, and the symmetry should not be over-read.
We do not claim that the mathematics of the release has been understood or reviewed, only that this objection to it is invalid; and what the objection uncovered about the distance between the formal proof and its own vocabulary remains worth the attention of anyone who sets out to explain the proof.

\section{Addendum (24 November 2026): the minimal example}
\label{sec:addendum}

After this note was completed, the author of \N\ posted a Lean file of about forty lines~\cite{nielsenpost}, described in its own comments as a minimal example with no \lean{sorry} and no custom axiom, and presented it as showing how the release compiles.
This addendum examines that file. The result can be stated at once.

\begin{verdictbox}
The example does not model the release. Its theorem \emph{assumes} the Riemann Hypothesis and hands the assumption back; the theorem of the release has no hypothesis, in the source and in the kernel export.
With the hypothesis removed the example does not compile, in any of the ways that were tried.
The objection therefore fails, and for the reason already given in Proposition~\ref{prop:vac}.
\end{verdictbox}

One further kind of evidence is used below: \textsf{L}, the output of the Lean compiler on a small file. The files and outputs are in \path{evidence/minimal-example} in the repository.

\subsection{The example}

The file defines a predicate \lean{anyGeometry} that is constantly \lean{True}, a number \lean{targetFormula}, a structure with a field \lean{height} and a field saying that \lean{height} equals \lean{targetFormula},
an instance in which that equation holds by \lean{rfl}, and a bundle of the two.
These are miniatures of objects that are in the release and on the route (Section~\ref{sec:facts}); agent~K matched each of them to its counterpart there.
The file ends with the following theorem and with the command \texttt{\#print axioms claimed\_rh}.
\begin{Verbatim}[fontsize=\small,frame=single,rulecolor=\color{gray!50}]
theorem claimed_rh (hRH : RiemannHypothesis) :
    (∃ b : Bundle, anyGeometry 0 ∧ b.model.height = targetFormula) ∧ RiemannHypothesis := by
  let b := buildBundle
  exact ⟨⟨b, valid_bundle_of_manufactured b⟩, hRH⟩
\end{Verbatim}
\noindent The comment above it in the file says that when the Riemann Hypothesis is combined with the certificate of the bundle, the certificate adds no independent evidence. That is true, and it is what the file shows.

\subsection{What the example states}

We elaborated the file with Lean~4.33.1 against the build of Mathlib that the release pins. The post displays the text flush left, which Lean does not parse; we restored the indentation and changed nothing else. Table~\ref{tab:minimal} collects the outcomes.

\begin{table}[ht]
\centering\small
\renewcommand{\arraystretch}{1.22}
\begin{tabularx}{\textwidth}{@{}>{\raggedright\arraybackslash}p{4.3cm}>{\raggedright\arraybackslash}X>{\raggedright\arraybackslash}p{1.6cm}@{}}
\toprule
Check & Outcome & Evidence \\
\midrule
The type of \lean{claimed_rh} & \texttt{RiemannHypothesis → (∃ b, …) ∧ RiemannHypothesis}: an implication whose hypothesis is its own conclusion. & \textsf{L}, \textsf{G-I} \\
Its proof term & \texttt{fun hRH => … ⟨Exists.intro b …, hRH⟩}: the second component is the bound variable. Nothing from the bundle and nothing from Mathlib enters it. & \textsf{L}, \textsf{G-I} \\
The same text with an arbitrary proposition \texttt{P}, or with \lean{False}, in place of \lean{RiemannHypothesis} & Compiles unchanged, with type \texttt{False → (…) ∧ False} in the second case. The text is not about the Riemann Hypothesis. & \textsf{L}, \textsf{G-I} \\
The three axioms reported by \texttt{\#print axioms} & They enter through the \emph{statement}, in which the definition of the zeta function occurs. With \lean{False} in its place the report is ``does not depend on any axioms''. & \textsf{G-I}, \textsf{L} \\
The bundle deleted altogether & \texttt{theorem t (hRH : RiemannHypothesis) : RiemannHypothesis := hRH} compiles, with the same axiom report. The bundle carries nothing. & \textsf{G-I}, \textsf{L} \\
The hypothesis left out: \texttt{theorem t : RiemannHypothesis := (claimed\_rh \_).2} & Error: the compiler cannot supply the argument \lean{hRH}, and shows the open goal \lean{RiemannHypothesis}. & \textsf{L}, \textsf{G-I} \\
Six further attempts to do without the hypothesis while keeping every vacuous ingredient & All six fail: extracting the conclusion from the bundle; a bundle field of type \lean{RiemannHypothesis}; \lean{rfl}, \lean{trivial}, \lean{decide}; a consumer quantified over all propositions; fields that mention the conclusion vacuously; redefining \lean{RiemannHypothesis}. & \textsf{G-K}, \textsf{L} \\
\bottomrule
\end{tabularx}
\caption{What the Lean compiler says about the example. \textsf{L}: compiled by us; \textsf{G-x}: established by agent~x and, where \textsf{L} follows, recompiled by us with the same outcome.}
\label{tab:minimal}
\end{table}

The example is thus an instance of Proposition~\ref{prop:vac}: a statement $H \to P$ whose proof says nothing about $P$. Here $H$ is the Riemann Hypothesis itself and $P$ contains it as a conjunct.
The certificate about the bundle is true and idle; the conclusion is the hypothesis.
What makes the file compile is the binder \texttt{(hRH : RiemannHypothesis)}. The predicate that is \lean{True} and the equation that holds by \lean{rfl} contribute nothing to that.

\subsection{What the release states}

The entry point of the release is the file \path{src/Main.lean}, and its first theorem reads, in full,
\begin{Verbatim}[fontsize=\small,frame=single,rulecolor=\color{gray!50}]
theorem riemannHypothesis : RiemannHypothesis := rhOfAll
\end{Verbatim}
\noindent There is no binder, and the file contains no \lean{variable} or \lean{section} that could add one (\textsf{S}, \textsf{G-I}, \textsf{G-J}).
The source is not what was checked, so we looked at the kernel export as well. In the 61.57\,GB proof export the declaration is recorded as
\begin{Verbatim}[fontsize=\scriptsize,frame=single,rulecolor=\color{gray!50},breaklines=true,breakanywhere=true]
{"thm":{"all":[14065054],"levelParams":[],"name":14065054,"type":6490658,"value":1065212062}}
{"const":{"name":326825,"us":[]},"ie":6490658}
{"in":326825,"str":{"pre":0,"str":"RiemannHypothesis"}}
\end{Verbatim}
\noindent that is, name 14065054 (\lean{riemannHypothesis}) is a theorem whose type is expression 6490658, and that expression is the bare constant \lean{RiemannHypothesis} (\textsf{K}).
A hypothesis would make the type a \texttt{forallE} record. It does so for the lemma \lean{Formalism.VW.rh_of_every} one step below, whose type is recorded as
\texttt{forallE} with domain \lean{UpperTheory.EveryCharacter} and body 6490658 (found by agent~M; confirmed by us, \textsf{K}).
The challenge export, which we rebuilt from Mathlib alone (Section~\ref{sec:rebuild}), records the same constant as the type of the same name (\textsf{K}; independently \textsf{G-L} and \textsf{G-M}).
Agent~M also found that the two exports agree byte for byte up to and including the definition of \lean{RiemannHypothesis}, their first 6,426,463 lines; we recomputed the two hashes and they coincide (\textsf{K}).
Its full pass over the proof export found three \texttt{axiom} records, the standard ones, and no occurrence of \lean{sorryAx}, in agreement with Section~\ref{sec:defects}.

This is not in dispute with \N. Its abstract says of the top-level theorem that it inhabits Mathlib's \lean{RiemannHypothesis}; that is observation (A1) of Section~\ref{sec:argument}.
A theorem of type \texttt{RiemannHypothesis → …} does not inhabit \lean{RiemannHypothesis}. The example therefore departs from the description of the release that \N\ itself gives.

\subsection{Is the Riemann Hypothesis assumed anywhere on the route?}

The example would be a fair model if the release, somewhere below its final theorem, assumed what it concludes. We looked for that in four ways.

\begin{enumerate}
\item \textbf{In the dependency graph.} Of the \Ndecl\ declarations of the proof export, exactly four refer to the constant \lean{RiemannHypothesis} in their type:
\lean{Formalism.VW.rh_of_right}, \lean{Formalism.VW.rh_of_every}, \lean{rhOfAll} and \lean{riemannHypothesis} (\textsf{K}). In each of the four it is the conclusion (\textsf{S}). None takes it as a hypothesis.
\item \textbf{In the source.} Agent~J classified all 51 occurrences of the identifier in the repository of the release: 47 as the conclusion of a statement, one inside a proof, three in documentation, none as a hypothesis (\textsf{G-J}).
The source contains more declarations than the export, which holds only what the final theorems use.
\item \textbf{In unfolded form.} An assumption need not use the identifier. Lemmas on the route do have hypotheses that say what the Riemann Hypothesis, or its generalization to Dirichlet characters, says:
\lean{rh_of_right} assumes that the zeta function has no zero with real part greater than $1/2$ other than at $s = 1$, and \lean{rh_of_every} assumes the corresponding statement for all Dirichlet $L$-functions.
Agent~J followed the route for twelve levels, listed seven hypotheses of this kind, and for each named the theorem that supplies it; agent~O, reading independently, found the same for the part of the route it covered (\textsf{G-J}, \textsf{G-O}).
This is what the last steps of any proof look like, and it is where the example and the release part ways.
In the release the chain is closed: \texttt{rhOfAll := rh\_of\_every allChi}, where \lean{allChi} is a theorem without hypotheses whose type is the constant \lean{UpperTheory.EveryCharacter} (\textsf{S}; in the export, \textsf{G-M}),
and the dependency cone of \lean{riemannHypothesis} contains \NconeProjThm\ theorems of the release (Appendix~\ref{app:data}).
In the example nobody supplies \lean{hRH}.
\item \textbf{By compiling the last steps.} The two lemmas just named use only Mathlib; the first is a proof of some sixty lines from the functional equation.
Copied verbatim into a file of their own they compile, with the three standard axioms and the types
\texttt{(∀ s, 1/2 < s.re → s ≠ 1 → riemannZeta s ≠ 0) → RiemannHypothesis} and \texttt{UpperTheory.EveryCharacter → RiemannHypothesis} (\textsf{L}; independently \textsf{G-O}).
The definition of \lean{UpperTheory.EveryCharacter} mentions only Mathlib's \lean{DirichletCharacter.LFunction} (\textsf{S}, \textsf{G-O}).
A last line that asks for \lean{RiemannHypothesis} without a proof of \lean{UpperTheory.EveryCharacter} fails in the same way as the example without \lean{hRH}.
\end{enumerate}

None of this requires reading every lemma of the release. A closed term that a kernel accepts with three axioms cannot contain an undischarged hypothesis; items 1--4 show where the hypotheses of the last steps are discharged, and the replay of Section~\ref{sec:conleche} covers the rest.

\subsection{The converse experiment}

Agent~K also ran the experiment in the other direction. It kept every vacuous ingredient of the example and passed the bundle, with two proofs of \lean{True}, to a lemma stated for all bundles whose conclusion is proved independently: the irrationality of $\sqrt{2}$, from Mathlib.
The resulting theorem has no hypothesis, has \lean{buildBundle} and \lean{anyGeometry} in its proof term, and is valid (\textsf{G-K}, \textsf{L}). Agent~L did the same with Euclid's theorem on the infinitude of primes (\textsf{G-L}, \textsf{L}).
Together with the six failures of Table~\ref{tab:minimal} this is the replacement test of Section~\ref{sec:remove} in miniature. Vacuous objects in the dependency cone of a closed theorem do not invalidate it, and without a proof of the conclusion they cannot produce it.

\subsection{What the example does show}

Two things, and both were already common ground.

\begin{enumerate}
\item \textbf{A vacuous certificate adds no evidence to a statement to which it is attached.} Section~\ref{sec:credit} says the same of the geometric layer of the release.
Alpöge has since written that the source contains much that is unnecessary and that criticism of it is reasonable~\cite{alpogepost}. That concession concerns the same layer. Like the layer, it bears on the state of the source and not on what the checked term proves.
\item \textbf{A clean axiom report does not show that a theorem proves what its name suggests.} Agent~L named this as the strongest point in favour of the example.
\lean{claimed_rh} is a valid theorem with no \lean{sorry} and only the standard axioms. It is simply not the Riemann Hypothesis. One has to read the statement.
That is exit~(ii) of Proposition~\ref{prop:exits}, and the mechanical check for it is the comparison of the statement with a reference built from Mathlib alone (Sections~\ref{sec:rebuild} and~\ref{sec:conleche}).
The release passes that check. The example would not: its type is not the constant \lean{RiemannHypothesis}.
\end{enumerate}

\subsection{Independent verification of the addendum}
\label{sec:agents3}

As in the main audit, the checking was shared with a second line of verification, and it corrected the first.
Six further agents were run with Gemini~4 Neon through \texttt{agy}, each from an empty context and on one sub-question. They were given the example and the claim made for it in neutral terms, were told not to read our earlier reports or each other's directories,
and, unlike the agents of Section~\ref{sec:agents}, could compile small Lean files against Mathlib. They could not build the release. (There is no agent~N, to avoid confusion with \N.)

\begin{table}[ht]
\centering\small
\renewcommand{\arraystretch}{1.22}
\begin{tabularx}{\textwidth}{@{}c>{\raggedright\arraybackslash}p{3.0cm}>{\raggedright\arraybackslash}X@{}}
\toprule
Agent & Sub-question & What it established \\
\midrule
I & what the example states and proves & The type, the proof term, the substitution of \texttt{P} and \lean{False}, the deletion of the bundle.
\textbf{Caught a vacuous check in our draft}: we had tested that the second component of \lean{claimed_rh} equals \lean{hRH} by \lean{rfl}, which holds for any two proofs of a proposition; it printed the term instead.
\textbf{Observed that the three reported axioms come from the statement, not from the proof.} \\
J & is the conclusion assumed on the route? & \textbf{Classified all 51 occurrences} of the identifier; \textbf{listed the seven hypotheses of this shape on the route and the theorem that discharges each}; found no construct that could add a hidden hypothesis in the 53 files of the chain. \\
K & is the example a faithful model? & Matched each ingredient to its counterpart in the release and \textbf{found none for the hypothesis} \lean{hRH} or for the final conjunction; \textbf{the six failed attempts}; \textbf{the converse experiment}. \\
L & logic of the inference & Conditions under which the inference from the example would be valid; verified the type in the challenge export; verdict: invalid. \textbf{The point in favour of the example} recorded above. \\
M & the kernel exports & Resolved the records of \lean{riemannHypothesis}, \lean{rhOfAll}, \lean{rh_of_every} and \lean{allChi} in the proof export, with the same result as ours;
\textbf{found the common prefix of the two exports}; a full pass over the 61.57\,GB file for axiom records. \\
O & where the conclusion comes from & Compiled the last steps on its own; unfolded \lean{UpperTheory.EveryCharacter} to Mathlib terms; walked the route, confirmed that the vacuous objects are on it and that their consumers are universally quantified, as in Section~\ref{sec:remove}. \\
\bottomrule
\end{tabularx}
\caption{The six agents of the addendum (Gemini~4 Neon). Bold marks what originated with the agent.}
\label{tab:agents3}
\end{table}

All six reached the assessment stated at the head of this addendum. The qualifications of Section~\ref{sec:agents} apply again: the agents read the same artefacts as we did, and the prompts were ours.
We recompiled all 23 Lean test files written by the agents and obtained in every case the outcome they report; the list is in the repository.
Agents~L and~O also recorded further definitions equal to \lean{True} on or near the route, among them a slot named \lean{UnusedSlot72}. They add to the list of Section~\ref{sec:facts} and do not change its assessment.

\subsection{Limits of the addendum}

The post is informal, and we have addressed only the file and the claim that it shows how the release compiles.
The limits of Section~\ref{sec:limits} apply unchanged: the mathematics was not read, and the build from source to export was not reproduced here.
Agent~J read the route to a depth of twelve levels and not to its leaves, so item~3 above is a sample; the absence of an assumed conclusion rests on items~1 and~2 and, decisively, on the fact that the checked term is closed.


\section*{Contributions, tools and acknowledgements}
\addcontentsline{toc}{section}{Contributions, tools and acknowledgements}
\label{sec:contrib}

This audit is the product of a person and two AI systems, and the division of labour should be stated plainly.

The author posed the question, supplied the paper under review, set the scope and the standard of evidence, directed the work and is responsible for the text and its conclusions.

\emph{Claude Code} (Anthropic) carried out the main line of analysis under the author's direction: it wrote the programs \texttt{conescan}, \texttt{depgraph} and \texttt{depquery}, ran them and the \texttt{con-leche} replay,
rebuilt the challenge export, selected the source excerpts and drafted this text. For the addendum it compiled the example and the last steps of the release, and looked up the records in the two exports.

\emph{Gemini~4 Neon} (Google), run as fourteen independent agents through the \texttt{agy} tool, eight for the main audit and six for the addendum, performed the independent verification described in Sections~\ref{sec:agents} and~\ref{sec:agents3}.
Several findings used in this note originated with those agents and not with the main analysis: the location of the degenerate and junk models, the observation that the custom tactics call for an independent kernel replay,
the substantive nature of the copied lattices, the two mis-cited sites, the link between \lean{Bk} and the period pairing, and the minimal counterexample to the deletion test;
and, in the addendum, the six failed attempts to remove the hypothesis of the example, the converse experiment, the classification of every occurrence of the conclusion in the source, and the common prefix of the two exports.
The main analysis, for its part, made an error in the first version of this audit, an undercount that came to light when the dependency graph was rebuilt with a second program (Appendix~\ref{app:erratum}),
and a vacuous check in its draft of the addendum, which agent~I caught. The results are stronger for having been obtained twice, by different systems.

Neither AI system is an author, and neither can take responsibility for the conclusions; the author does.
The author thanks J.~L.~Nielsen for a paper whose careful reading of the source made this audit possible, and the maintainers of the verification page and of the mirrors for making the artefacts available.

\begin{thebibliography}{9}
\bibitem[N]{nielsen} J.~L.~Nielsen, \emph{The Invalidity of Alpöge's Proof of the Riemann Hypothesis: Vacuously true geometric and lattice witnesses in the Lean formalization},
Center for Topological Physics, Lawrence, Kansas; title page dated 20 November 2026, 17~pages. PhilArchive record \url{https://philarchive.org/rec/NIEIOT-2}.
\bibitem{nielsenpost} J.~L.~Nielsen, post on X containing a minimal Lean example, 24 November 2026. \url{https://x.com/QualiaQuanta/status/2125008805310222352}
\bibitem{alpogepost} L.~Alpöge, post on X, 24 November 2026. \url{https://x.com/__alpoge__/status/2124987173439209561}
\bibitem{alpoge} L.~Alpöge, \emph{Remark}, 20 November 2026. \url{https://alpo.ge/remark.pdf}
\bibitem{community} Lean and Mathlib communities, statement on the formalization, 21 November 2026. \url{https://leanprover-community.github.io/blog/posts/riemann-hypothesis/}
\bibitem{mit} A.~V.~Sutherland, verification files for the release (proof export, challenge export, checker logs). \url{https://math.mit.edu/~drew/grh/}
\bibitem{mirror} Mirrors of the release: \url{https://github.com/the-omega-institute/riemann-hypothesis} (commit \texttt{c7d25b3b}, used here) and \url{https://github.com/Deicyde/riemann-hypothesis} (commit \texttt{712a6aef}, cited in \N).
\bibitem{conleche} \texttt{con-leche}, an independent checker for Lean~4 exports. \url{https://github.com/leanprover/con-leche}
\bibitem{jaswal} M.~Jaswal, \emph{The Riemann hypothesis for Dirichlet $L$-functions [after Levent Alpöge]}, version~1, 21 November 2026; cited in \N. Not examined in this note.
\bibitem{repo} Repository of this audit: programs, outputs, logs, agent prompts and reports. \url{https://github.com/ArtificialZeng/Riemann-hypothesis-Auditing}
\end{thebibliography}


\appendix
\section{Kernel-level data}
\label{app:data}

All numbers in this appendix are produced by \path{evidence/depgraph/depquery.py} from the dependency graph of the export and inserted into the text by \path{paper/make_tables.py}.

\subsection{The export}

\begin{center}\small
\begin{tabular}{@{}lr@{\qquad}lr@{}}
\toprule
\multicolumn{2}{c}{\textbf{Size}} & \multicolumn{2}{c}{\textbf{Declaration records}}\\
\midrule
lines & \Nlines & theorems & 1,246,110 \\
bytes & 61,570,595,788 & definitions & 214,241 \\
expression nodes & \Nexprs & inductive blocks & 7,969 \\
type-dependency edges & \NedgesT & opaque constants & 91 \\
value-dependency edges & \NedgesV & quotient primitives & 4 \\
library declarations & \Nlib & axioms & 3 \\
declarations of the release & \Nproj & \textbf{total} & \textbf{\Ndecl} \\
\bottomrule
\end{tabular}
\end{center}

The dependency cone of \lean{riemannHypothesis} alone has \NconeMain\ declarations (\NconeProj\ of the release, of which \NconeProjThm\ are theorems); the few remaining records belong to the statements of the other four theorems.
Its statement cone has \NstmtMain\ declarations, \NstmtMainProj\ of them from the release.

\subsection{Dependency cones along the route}

\begin{center}\footnotesize
\begin{tabular}{@{}lrrrc@{}}
\toprule
Declaration & cone & release & theorems & flagged? \\
\midrule
\tableinput{tables/route.tex}
\bottomrule
\end{tabular}
\end{center}
\noindent The columns give the size of the dependency cone, the number of its declarations that belong to the release, the number of those that are theorems, and whether the cone contains a flagged object, meaning any of \lean{anyGeometry}, \lean{geoDegenerate}, \lean{algDegenerate}, \lean{padDegenerate}, \lean{mdlJunk}, \lean{cxJunk}, \lean{divXFl}, \lean{geoOver_stages}, \lean{geoExample}, \lean{pointGeo}, \lean{x0}, \lean{provisionalDivX}, \lean{graft}.
The following names, all mentioned in \N\ or in earlier accounts, do not occur in the export: \tableinput{tables/absent.tex}.

\subsection{Dependents of individual objects}

\begin{center}\small
\begin{tabular}{@{}lrc@{}}
\toprule
Object & declarations depending on it & final theorem among them \\
\midrule
\tableinput{tables/dependents.tex}
\bottomrule
\end{tabular}
\end{center}
\noindent The two structures whose names end in \texttt{Placeholder} are discussed in Appendix~\ref{app:erratum}. Excluding them, the union of the sets above has \Nflagged\ elements (\NflaggedThm\ theorems).
For comparison, \NifaceDeps\ declarations depend on the interface type \lean{DV.Geo} and \NdivXdeps\ on the record \lean{DivX}.

\subsection{Direct references to selected names}

\begin{center}\small
\begin{tabular}{@{}lc>{\raggedright\arraybackslash}p{8.2cm}@{}}
\toprule
Name & referrers outside its own block & which \\
\midrule
\tableinput{tables/watch.tex}
\bottomrule
\end{tabular}
\end{center}

\subsection{The \NanyGeo\ declarations that depend on \texttt{anyGeometry}}

{\scriptsize
\begin{longtable}{@{}rl>{\raggedright\arraybackslash}p{13.2cm}@{}}
\toprule
\# & kind & declaration \\
\midrule
\endhead
\bottomrule
\endfoot
\tableinput{tables/anygeo.tex}
\end{longtable}
}

\section{Correction to the earlier version of this audit}
\label{app:erratum}

The earlier version (in Chinese, dated 23 November 2026) reported some counts obtained with \texttt{conescan} that were too low. The cause and the extent are as follows.

\texttt{conescan} assigns a taint mark to an expression node when the node is created. The exporter emits the expressions of an inductive declaration, including the node for the type's own name, \emph{before} the declaration record.
When an inductive type depended on a seed, that early node had therefore been created unmarked, and later declarations that referred to the type through the same shared node were not counted.
Propagation through definitions and theorems was not affected, and neither was the detection of direct references to a seed.

\begin{itemize}
\item \textbf{Unaffected} (re-derived with \texttt{depgraph} and identical): \NanyGeo\ dependents of \lean{anyGeometry}; \NgeoDeg\ of \lean{geoDegenerate}; the counts for \lean{algDegenerate}, \lean{padDegenerate}, \lean{mdlJunk}, \lean{cxJunk}, \lean{divXFl}, \lean{geoOver_stages}, \lean{geoExample};
the absence of \lean{plugAll} and of the other names listed above; the absence of any referrer of \lean{GeoCl.rec}; the two readers of \lean{OAt.geo}; the three axioms.
\item \textbf{Corrected.} The earlier text said that 736 declarations involve the type \lean{DV.Geo} and that 621 of them read one of its fields. The correct statements are those of Section~\ref{sec:degenerate}:
\NifaceDeps\ declarations depend on the type, \Nacc\ field accessors occur in the export, \NaccDirect\ declarations refer to an accessor directly, and \NaccThmProof\ of those are theorems whose proof does so.
The count given for \lean{GeoTiedP} (60) should read \NgeoTiedP; the counts given for field projections of \lean{DivX} and \lean{DivT} are superseded by the direct-reference data of Appendix~\ref{app:data}.
\item \textbf{Corrected.} Two structures with no fields, \lean{LPTiedPlaceholder} and \lean{NTiedGPlaceholder}, were reported with 117 and 82 dependents. The correct numbers are \NplaceLP\ and \NplaceNT:
they are parameters of record types that are used widely, including in the type of the bundle \lean{HypsAlt}. A structure without fields has exactly one inhabitant and carries no information, so nothing can be smuggled through it; \N\ does not mention them.
\end{itemize}

None of the corrections changes a conclusion, but the earlier wording understated how widely the abstract geometry \emph{interface} is threaded through the data types, and the distinction now drawn in Section~\ref{sec:degenerate} between the interface and its trivial instances is the accurate one.
The tool \texttt{depgraph} avoids the problem by resolving references when a declaration is complete, not when an expression is created.

\section{Hashes and versions}
\label{app:hashes}

\begin{center}\scriptsize
\renewcommand{\arraystretch}{1.25}
\begin{tabularx}{\textwidth}{@{}>{\raggedright\arraybackslash}p{4.3cm}>{\raggedright\arraybackslash}X@{}}
\toprule
Object & SHA-256 or commit \\
\midrule
\path{lean.tar.gz} (author's archive) & \texttt{9631b4f020199e580ba86178cb56c6e9b09f690c2369a80dc4334aa93f49c7fc} \\
\path{solution.ndjson.zst} & \texttt{61a8938ad5f5e57397c1c0e3b954916af546d9004ade01c3555744a681b05381} \\
\path{solution.ndjson} & \texttt{5ea74b912147c5c93cc2b6d70cc727cf908df2f62ac6a20480a3fa02f6be244e} \\
\path{challenge.ndjson} (published and rebuilt) & \texttt{bd484ac9886a96f26b1d9a7a59ebd45b7a70d102c015d4f31f9302bbdd6ab2b9} \\
The paper \N, version reviewed & \texttt{90b7c28e51e6cb8460480c7ecbb81162f1b910c24234f7664dcd1bafbc7ccf77} \\
Mirror used & \path{the-omega-institute/riemann-hypothesis}, commit \texttt{c7d25b3bfd7906388bb8eb0774e027a2b756a7b6} \\
Mathlib & \texttt{0df444a360eaa60ab8c11dca51a86af692955474} (official tag \texttt{v4.33.1}) \\
con-leche & \texttt{bc6676101c488e21ce71a317b33905534d497e14}, built with Lean 4.36.0 \\
lean4export & \texttt{15f6055e299ad5b89345e533cc2192f4cc00f659} (tag \texttt{v4.33.0}), built with Lean 4.33.1 \\
Dependency graph \path{edges.bin} & \texttt{\edgesHash} \\
\bottomrule
\end{tabularx}
\end{center}

\section{Reproduction}
\label{app:repro}

The file \path{evidence/REPRODUCE.md} in the repository gives the commands. In outline (about 80\,GB of disk; the checker run needs about 48\,GB of memory):
\begin{enumerate}
\item download the fourteen parts of the archive and the two exports, and compare the hashes above;
\item \texttt{cargo build --release} in \path{evidence/conescan}, then \texttt{depgraph OUT watch\_names.txt < solution.ndjson} (one pass, about eight minutes, 17\,GB of memory);
\item dump the Mathlib constant names with \path{evidence/depgraph/DumpNames.lean}, run \path{depquery.py}, then \path{paper/make_tables.py};
\item build \texttt{con-leche} and run it with \texttt{--verified --compare};
\item rebuild the challenge export with \path{evidence/challenge-rebuild/run2.sh} and compare it byte for byte;
\item compile the two files in \path{evidence/toy} with any recent Lean~4;
\item for the addendum, compile the files in \path{evidence/minimal-example} inside the project of step~5 and run \path{rh_referrers.py} on the graph of step~2.
\end{enumerate}


\end{document}
